\documentclass[12pt,a4paper]{article}
\usepackage[utf8]{inputenc}
\usepackage[T1]{fontenc}
\usepackage[ngerman,latin]{babel}
\usepackage{amsmath,amssymb,amsthm}
\usepackage{geometry}
\usepackage{hyperref}
\usepackage{enumitem}
\usepackage{mathrsfs}
\usepackage{longtable}
\usepackage{booktabs}
\usepackage{microtype}
\usepackage{array}
\usepackage{colortbl}
\usepackage{graphicx}
\DeclareMathOperator{\cotang}{cotang}
\geometry{margin=2.5cm}

\theoremstyle{plain}
\newtheorem{satz}{Satz}[section]
\newtheorem{theorem}{Theorem}[section]
\newtheorem{lemma}{Lemma}[section]
\newtheorem{hilfssatz}{Hilfssatz}[section]
\newtheorem{korollar}{Korollar}[section]

\theoremstyle{definition}
\newtheorem{definition}{Definition}[section]

\theoremstyle{remark}
\newtheorem{bemerkung}{Bemerkung}[section]
\newtheorem{beispiel}{Beispiel}[section]

\title{Band 1 (Alternate), Teil 8}
\author{Carl Friedrich Gau\ss}
\date{}

\begin{document}
\maketitle

\section*{Bemerkungen. Dreieckskranz um Oldenburg.}
\addcontentsline{toc}{section}{Bemerkungen. Dreieckskranz um Oldenburg.}

Die Seitengleichung $1^{\circ}$ auf S.~336 ist die Bedingung daf"ur, dass der Werth der Seite $s_r$ wieder erhalten wird, wenn durch das Kranzsystem hindurch gerechnet wird. Es ist

\begin{equation*}
s_r^{\circ} = s_1^{\circ} \frac{\sin a' \cdot \sin b' \cdot \sin c'}{\sin a'' \cdot \sin b'' \cdot \sin c''} \cdots
\end{equation*}

also
\begin{multline*}
\delta s_r = -\delta s_1 = \cotang a' \cdot \delta a' - \cotang a'' \cdot \delta a'' + \cdots \\
+ \cotang b' \cdot \delta b' - \cotang b'' \cdot \delta b'' + \cdots
\end{multline*}

F"ur den briggischen Logarithmus und wenn die Verbesserungen $\delta a'$, $\delta a''$, etc.\ in Secunden erhalten werden sollen, ist die linke Seite dieser Gleichung mit $\frac{206265}{\rho''}$ zu multipliciren.

In den Formeln 1.\ und 2.\ S.~336 und 337, steht im Original
\[\frac{e^{1-2\lambda}}{1-2\lambda}\]
an Stelle von $-\frac{e^{1-2\lambda}}{1-2\lambda}$ und $\frac{a^{1-2\lambda}}{1-2\lambda}$.

Im Jahre 1834 hat Gauss nach dem \emph{Generalbericht "uber die mitteleurop"aische Gradmessung f"ur das Jahr 1868}, S.~22/23, dem oldenburgschen Vermessungsdirector v.~Schrenck die 21 sph"aroidischen, auf $160''$ Excess abgestimmten Dreiecke des Kranzes um Oldenburg mitgetheilt. In einem Schreiben vom 3.~August 1835 sagt Gauss dazu (nach dem angegebenen Generalbericht, S.~23/24):

\begin{quote}
>>Was nun aber die relative Genauigkeit der Dreiecke betrifft, so sind Ew.\ Hochw.\ im Irrthum, wenn Sie die Dreiecke 2--6 (in der Figur auf S.~332; sind dies die Dreiecke $g$, $p$, $o$, $n$, $m$, $l$, $k$) den "ubrigen 9--21 und 1 (in der Figur $c$ bis $a$, $v$ bis $r$) entgegenstellen. Der Gegensatz soll vielmehr so sein: 3--16 (in der Figur $p$ bis $b$) viel ungenauer, als 17--21 ($a$, $v$, $u$, $t$, $s$) und 1 ($r$) und 2 ($q$). Die letztern 7 Dreiecke habe ich selbst gemessen mit 12-z"olligen Theodolithen, gr"osster Sorgfalt, Heliotroplicht ohne Ausnahme die Zielpunkte bildend, und unter m"oglichster Sorge f"ur Festigkeit der Standpunkte, wovon 3 zu ebener Erde. Dagegen sind die 14 andern Dreiecke zu anderm Zwecke, mit schw"acherem Instrument (8-z"olligem Theodolith), viel geringerm Zeitaufwand, ohne Anwendung von Heliotroplicht und mitunter auf sehr ung"unstigen Standpunkten gemessen, wie z.~B.\ die Th"urme von Twistringen und Asendorf und vielleicht auch einige der andern Th"urme. Indem ich daher die 7 ersten Dreiecke f"ur so scharf gemessen halte, wie das der Zustand der Kunst nur verstattet, w"urde ich die Genauigkeit der 14 "ubrigen nur so gross, oder ihr Gewicht nur $\frac{4}{7}$ so gross ansetzen.<<
\end{quote}

Wenn man die Winkelsumme des Siebenecks 4, 6, 10, 12, 16, 19, 21 aus den von Gauss an v.~Schrenck mitgetheilten Winkeln bildet, so betr"agt (Generalbericht, S.~24) der Schlussfehler $1\frac{4}{9}''$.

\begin{flushright}
\textsc{Kriiger.}
\end{flushright}

\newpage
\section*{Briefwechsel zur Hannoverschen Triangulation.}
\addcontentsline{toc}{section}{Briefwechsel zur Hannoverschen Triangulation.}

\subsection*{[1.]}
\addcontentsline{toc}{subsection}{Gauss an Schumacher}

\textbf{Gauss an Schumacher.} G"ottingen, 5.~Julius 1816.

\medskip
\noindent Vor allen Dingen meinen herzlichen Gl"uckwunsch zu der herrlichen grossen Unternehmung, welche Sie mir in Ihrem letzten Briefe ank"undigen. Diese Gradmessung in den k.\ d"anischen Staaten wird uns, an sich schon, "uber die Gestalt der Erde sch"one Aufschl"usse geben. Ich zweifle indessen gar nicht, dass es in Zukunft m"oglich zu machen sein wird, Ihre Messungen durch das K"onigreich Hannover s"udlich fortzusetzen. In diesem Augenblicke kann ich zwar einen solchen Wunsch in Hannover noch nicht in Anregung bringen, da erst die Astronomie selbst noch so grosser Unterst"utzung bedarf; allein ich bin "uberzeugt, dass demn"achst unsere Regierung, die auch die Wissenschaften gern unterst"utzt, dem glorreichen Beispiele Ihres trefflichen K"onigs folgen werde. Wir w"urden dann schon einen respectabeln Meridianbogen von $6\frac{1}{2}$ Grad haben, und leicht w"urden sich dann auch noch diese Operationen mit den bayerischen Dreiecken in Verbindung setzen lassen. Letztere sind gewiss mit gr"osster Sorgfalt gemessen, und es ist zu beklagen, dass sie der Publicit"at entzogen werden.

"Uber die Art, die gemessenen Dreiecke im Calcul zu behandeln, habe ich mir eine eigene Methode entworfen, die aber f"ur einen Brief viel zu weitl"aufig sein w"urde. In Zukunft, falls ich bis dahin, wo Sie Ihre Dreiecke gemessen haben, sie nicht schon "offentlich bekannt gemacht haben sollte, werde ich mit Ihnen dar"uber umst"andlich conferiren; ja ich erbiete mich, die Berechnung der Hauptdreiecke selbst auf mich zu nehmen.

\medskip
\noindent \textsc{IX.}~44

\newpage
\noindent Bei dem zweiten Theile Ihrer Unternehmung, der Messung des L"angengrades, habe ich nur einen kleinen Zweifel. Ich meinte nemlich, dass die L"ander der d"anischen Monarchie eher flach zu nennen sind, wenigstens keine hohe Berge haben. Ist diese Voraussetzung gegr"undet, und sind Sie dann dadurch gen"othigt, zur Bestimmung des astronomischen L"angenunterschiedes einen Zwischenpunkt oder gar mehrere zu nehmen, so wird jener Bestimmung, auch wenn sie noch so geschickte Gehilfen und H"ulfsmittel haben, doch immer eine kleine Ungewissheit ankleben.

Einen grossen Vortheil haben Sie in dem Umstande, dass D"anemark schon einmal trigonometrisch vermessen ist; ich meine nat"urlich nicht in den gemessenen Winkeln selbst, die weit davon entfernt sind, sich zu einer Gradmessung zu qualificiren, sondern weil jene Operationen Ihnen das Ausw"ahlen der Stationspunkte ungemein erleichtern werden. Dies Aufsuchen w"urde mir bei einer "ahnlichen Arbeit gerade das Unangenehmste sein, weil dabei so viele Zeit umsonst verloren wird. Ich habe mir viele M"uhe gegeben (in "ahnlichen R"ucksichten auf k"unftige Operationen), die von Eimbeck im Hannoverschen gemessenen Winkel zu erhalten, aber ohne Erfolg.

\medskip
\noindent Mir war eine interessante Aufgabe eingefallen, nemlich:
\begin{quote}
>>allgemein eine gegebene Fl"ache so auf einer andern (gegebenen) zu projiciren (abzubilden), dass das Bild dem Original in den kleinsten Theilen "ahnlich werde.<<
\end{quote}
Ein specieller Fall ist, wenn die erste Fl"ache eine Kugel, die zweite eine Ebene ist. Hier sind die stereographische und die merkatorsche Projection particul"are Aufl"osungen. Man will aber die allgemeine Aufl"osung, worunter alle particul"aren begriffen sind, f"ur jede Art von Fl"achen.

Es soll hier"uber in dem \emph{Journal philomathique} bereits von Mong\'e und Poinsot gearbeitet sein (wie Burckhardt an Lindenau geschrieben hat), allein da ich nicht genau weiss, wo, so habe ich noch nicht nachsuchen k"onnen, und weiss daher nicht, ob jener Herren Aufl"osungen ganz meiner Idee entsprechen und die Sache ersch"opfen. Im entgegengesetzten Fall schiene mir dies einmal eine schickliche Preisfrage f"ur eine Societ"at zu sein.\,\ldots\ldots\ldots

\newpage
\noindent \textbf{Gauss an Schumacher.} G"ottingen, 10.~September 1818.

\medskip
\noindent Ich eile Ihnen anzuzeigen, dass ich von unserm Minister Arnswald den Auftrag erhalten, die zur Verbindung einer hannoverschen Triangulirung mit der Ihrigen n"othigen Messungen in L"uneburg vorzunehmen und dazu das N"othige mit Ihnen zu verabreden. Er macht zugleich mir Hoffnung, dass demn"achst auch die Fortsetzung selbst wohl zu Stande kommen werde, und es freut mich, dass diese nun durch die in L"uneburg vorzunehmenden Operationen gesichert werden kann.\,\ldots\ldots\ldots

\medskip
\noindent \textbf{Gauss an Schumacher.} G"ottingen, 20.~Mai 1820.

\medskip
\noindent In dieser Ungewissheit\footnote{"Uber Schumachers Aufenthalt.} adressire ich diesen Brief nach Copenhagen und w"unsche sehnlich, dass er Sie treffen und bald treffen m"oge; er soll Ihnen nemlich die Nachricht anzeigen, dass in Folge eines Schreibens vom Grafen von M"unster aus London, als Antwort meines vor einem Jahre von Ihnen gef"alligst besorgten Briefes\footnote{Abgedruckt in Band IV, S.~482/483.},
\begin{quote}
>>der K"onig die Fortsetzung der Gradmessung durch das K"onigreich Hannover genehmigt hat<<.
\end{quote}

\medskip
\noindent \textbf{Gauss an Schumacher.} G"ottingen, 4.~M"arz 1821.

\medskip
\noindent Es war fr"uher meine Absicht, von Hamburg anzufangen und so von Norden nach S"uden zu messen, allein da ich leider so sehr durch die Schreibfaulheit und Unzuverl"assigkeit aller K"unstler, mit denen ich zu thun habe, hingehalten werde, und bis jetzt noch gar nichts von den n"othigen H"ulfsmitteln in H"anden habe, so w"urden die daraus erwachsenden Verlegenheiten noch viel gr"osser sein, wenn ich jenen Plan befolgte. Dieser und noch verschiedene andere wichtige Gr"unde n"othigen und bewegen mich zu dem umgekehrten Plan, von S"uden nach Norden zu messen\ldots\ldots\ldots

\newpage
\noindent \textbf{Gauss an Schumacher.} Wulfsode, 18.~September 1822.

\medskip
\noindent So viel eine vorl"aufige Inspection des Terrains urtheilen l"asst, w"urde es nicht unm"oglich sein, die ganze Linie von Breithorn bis Eschede [Scharnhorst] (11220 Meter lang) unmittelbar zu messen. Welch eine herrliche Basis w"are dies!\,\ldots\ldots\ldots

\medskip
\noindent \textbf{Gauss an Schumacher.} [G"ottingen, Januar 1825.]

\medskip
\noindent Wenn ich alle gr"ossern und kleinern Durchhaue aus den Jahren 1821--1824 zusammen z"ahle, von solchen, wo vielleicht ein Dutzend B"aume gef"allt sind, bis zu den gr"ossten, so m"ogen etwa 16 oder 17 Durchhaue vorgekommen sein. Der Allergr"osste, nach der Ausdehnung, war im Becklinger Holz unweit der Strasse von Bergen nach Soltau.\,\ldots\ldots\ldots

\medskip
\noindent \textbf{Gauss an Schumacher.} Dangast, 1 Stunde von Varel, 20.~Junius 1825.

\medskip
\noindent Ich selbst pflege [beim Heliotrop] durch die drei Fussspitzen einen Kreis zu beschreiben, dessen Centrum als Zielpunkt betrachtet wird. Bei meinen beiden neuesten Heliotropen ist noch das Centrum selbst durch eine Spitze bezeichnet, welches viel Bequemlichkeit verschafft\ldots\ldots\ldots

\medskip
\noindent \textbf{Gauss an Schumacher.} G"ottingen, 14.~Januar 1829.

\medskip
\noindent Da wir noch wenig dar"uber wissen, ob die Figur der Erde von der Figur des mittlern Ellipsoids in l"angern oder in k"urzern Undulationen abweicht, so bin ich jedenfalls der Meinung, dass man am besten thut, immer das mittlere Ellipsoid zum Grunde zu legen.

Wenn ich jedoch es nicht gerade unbedingt verwerfen will, einmal ein osculirendes Ellipsoid zum Grunde zu legen, so kann ich dies doch nur da f"ur zul"assig halten, wo man Mittel hat, ein osculirendes Ellipsoid zu bestimmen, d.\,i.\ ein solches, in dem die Kr"ummung sowohl im Sinne des Meridians, als in dem darauf senkrechten Sinn der wirklichen Gestalt so nahe wie m"oglich kommt. Dies ist aber bloss aus der Verbindung von Messungen in beiderlei Sinn zu erhalten (Breiten- und L"angen-Gradmessung).\,\ldots\ldots\ldots

\newpage
\noindent Wenn ein Meridian nicht wirklich elliptisch ist, so kann man allerdings eine Ellipse berechnen, die sich an zwei St"ucke eines Bogens anschliesst, und man mag dies meinethalben eine osculirende Ellipse nennen; allein durch Umdrehung dieser Ellipse um ihre kleine Axe\footnote{Von der man nur sagen kann, dass sie der Erdaxe parallel ist, ohne mit ihr zusammen zu fallen, ja von der sie in der Regel weit abstehen wird.} entsteht keine Fl"ache, die man osculirendes Ellipsoid nennen darf, oder mit andern Worten, zwischen dem wirklichen Werth des L"angengrads und demjenigen L"angengrad, den man auf dem durch Umdrehung jener osculirenden Ellipse [entstandenen Ellipsoid] berechnet, ist gar kein Zusammenhang. Durch Verwechselung beider setzt man sich den gr"ossten Fehlern aus.\,\ldots\ldots\ldots

\newpage
\subsection*{[2.]}
\addcontentsline{toc}{subsection}{Gauss an Bessel}

\noindent \textbf{Gauss an Bessel.} G"ottingen, 26.~December 1821.

\medskip
\noindent Der vorige Sommer ist gr"osstentheils mit Vorbereitungen zum Trianguliren und dem Trianguliren selbst zugebracht.\,\ldots\ Die Winkel in Sternwarte, Meridianzeichen, Hohehagen, Hils, Brocken sind gemessen; nur am letzten Punkte ist der Winkel zwischen Inselsberg und andern Punkten missgl"uckt, weil das unerh"ort schlechte Wetter w"ahrend der Zeit, wo dort der Heliotrop war, fast alles Beobachten untersagte. Nur einmal auf eine halbe Stunde konnte die Richtung k"ummerlich gemessen werden, was aber mit andern fr"uhern ebenso k"ummerlichen Messungen, wo einmal auf das Haus, das andere Mal auf Encke's Sextanten-Viceheliotrop 3 Minuten vor Sonnenuntergang pointirt wurde, nicht gut harmonirt.

Der Umstand, dass ich nur Einen Heliotrop zu meiner Disposition hatte, hielt ungemein auf. Bei meinem Aufenthalt auf dem Hils musste der Heliotrop successive von Lichtenberg zum Meridianzeichen, dem Brelingerberg und Deister "ubergehen; ebenso wie ich auf dem Brocken war, reiste jener von Lichtenberg zum Hils, Hohehagen, Inselsberg. Die wenigsten der Winkel haben also unmittelbar gemessen werden k"onnen; ich hatte auf diesen beiden Stationen mehrere andere an sich sichtbare Punkte ausgew"ahlt und maass allezeit, was sich eben messen liess. Ich sehe aus dem 4.\ Bande der \emph{Base du syst\`eme m\'etrique}, dass die Franzosen in Spanien es ebenso gemacht haben; sie haben aber die Messungen nicht richtig combinirt. Wenn ich im n"achsten Jahre meine Triangulirung fortsetze, wird es besser gehen, indem ich zwei Heliotrope mehr haben werde. Meinen zum Viceheliotrop eingerichteten Sextanten habe ich immer bei mir gef"uhrt und zu telegraphischen Zeichen gebraucht. Es geht damit ganz vortrefflich, nur dass, weil jener Viceheliotrop auf Winkelabst"ande unter etwa $138^{\circ}$ begrenzt ist, ich nicht immer zu jeder Tageszeit ihn brauchen konnte; k"unftiges Jahr w"urde ich zu diesem Zweck in der Regel einen der Heliotrope bei mir behalten.

Die gr"osste Entfernung, wo das Heliotroplicht mit blossen Augen gesehen, ist die vom Brocken zum Hohehagen, indem am letztern Orte meine auf dem Brocken gegebenen telegraphischen Zeichen so gesehen wurden; die Entfernung ist $9\frac{1}{2}$ geographische Meilen; es geh"oren dazu aber wohl g"unstige Umst"ande. Bei Entfernungen bis 6 Meilen habe ich (trotz des Verlustes an Licht beim Gebrauch der Lorgnette), wenn die Umst"ande nur leidlich g"unstig waren, das Heliotroplicht mit blossem Auge bequem gesehen, ebenso in den Vormittagsstunden das Heliotroplicht vom Hils zum Brocken ($7\frac{1}{2}$ Meilen). Die Spiegelfl"ache ist $2\frac{1}{2}$ Quadratzoll, bei den beiden neuen Heliotropen habe ich $6\frac{1}{4}$ Quadratzoll genommen; m"assigen kann man das Licht leicht, bei grossen Entfernungen kann es aber doch F"alle geben, wo das st"arkere Licht angenehm ist; so drang einen Tag das Licht vom Brelingerberg zum Hils nur kurze Zeit durch, obgleich (oder vielmehr richtiger weil) ein v"ollig wolkenloser Himmel war. An solchen Tagen ist in der Regel die Durchsichtigkeit der Atmosph"are am geringsten; ich erinnere mich eines solchen Tages, wo das Heliotroplicht vom Deister zwar recht gut zu sehen war, aber nicht benutzt werden konnte, weil keiner der "ubrigen Gegenst"ande, selbst nicht einmal das nur wenig "uber eine Meile entfernte Einbeck gesehen werden konnte.

Einen sehr wesentlichen Vorzug hat das Heliotroplicht vor jedem andern Signale, wor"uber ich oft artige Erfahrungen gemacht habe; nemlich das Heliotroplicht sieht man desto besser, je st"arker man vergr"ossert, irdische Signale hingegen (bei grossen Entfernungen) desto schlechter, denn bei letztern ist es vorz"uglich die Bl"asse, die das Sehen hindert; von Hannover habe ich z.~B.\ auf dem Brocken, obgleich ich den Platz nach H"ohe und Azimuth genau wusste, niemals eine Spur sehen k"onnen.

\newpage
\noindent An dem Tage, wo ich das k"ummerliche Heliotroplicht vom Inselsberg erhielt (bloss Schuld des Wetters, denn bei einigermassen g"unstiger Luft musste das Heliotroplicht von 24 Quadratzoll Fl"ache noch "uberreichlich hell sein), sah ich dies zwar im Fernrohr des Theodolithen noch ziemlich gut, konnte damit aber keine Spur vom Umriss des Berges erkennen; dagegen sah mein Gehilfe diesen Umriss ganz leidlich in einem an sich sehr elenden Fernrohr von schwacher Vergr"osserung, konnte damit aber das Heliotroplicht nur selten sehen. Die Sache erkl"art sich leicht; auch bei der st"arksten Vergr"osserung bleibt das Heliotroplicht ein Punkt und dessen Licht ist immer dasselbe, aber der Grund ist desto d"usterer, je st"arker man vergr"ossert etc.

Grosse Signale habe ich nur zwei gebaut, auf dem Hohehagen und Hils, und vermuthlich w"are auch dies unterblieben, wenn mein Heliotrop sechs Wochen fr"uher vollendet gewesen w"are. Das Hils-Signal projicirt sich vom Brocken aus gegen nahen dunkeln Hintergrund; w"ahrend des ganzen Monats, den ich auf dem Brocken zubrachte, habe ich jenes nur zweimal "uberhaupt sehen und nur auf etwa zwei Minuten so sehen k"onnen, dass ein Winkel sich h"atte messen lassen; das Hohehagen-Signal projicirt sich gegen die entferntern Casselschen Berge, war "ofters zu sehen und auch ein paar Mal zu beobachten, obwohl ich, h"atte ich nicht den Heliotrop dahin geschickt, auch meine Winkelmessung nicht voll bekommen h"atte. Ausser der schweren Sichtbarkeit, Kosten und Zeitaufwand haben die Signalth"urme noch eine sehr ungl"uckliche Seite, den Reiz, welchen sie dem rohen Muthwillen zur Zerst"orung darbieten. Leider ist mein Hohehagen-Signal seit kurzem fast ganz verw"ustet, und ich werde gl"ucklich sein, wenn ich nur den Punkt mit hinreichender Sch"arfe wiederfinden kann.

\medskip
\noindent \textbf{Gauss an Bessel.} G"ottingen, 15.~November 1822.

\medskip
\noindent Die ausserordentlichen Schwierigkeiten, ein Dreiecksnetz in der L"uneburger Heide zu fahren, kannte ich schon aus Eimbeck's Bericht, der es geradezu f"ur unm"oglich erkl"art und seine Dreiecke, um den s"udlichen Theil von Hannover mit Hamburg zu verbinden, "uber Bremen, die Weser herunter und so die Elbe wieder herauf gef"uhrt hat. Und doch hatte er grosse Vortheile vor mir voraus; er beobachtete durchaus oben in seinen Signalth"urmen, wo er sich viel leichter weitere Aussicht verschaffen konnte als ich, der "uberall zu ebener Erde gemessen hat; er brauchte beim Aushauen der Waldungen viel weniger auf Schonung R"ucksicht zu nehmen; "uberall freier Transport, freie Arbeiter, etc., w"ahrend ich alles ohne Ausnahme mit baarem Gelde bezahlen muss und oft die Kosten so sehr weit "uber meine Vermuthung hinausgehen sehe; er hatte eine ganze Brigade von Ingenieurs zu Gehilfen, etc. Diese Erw"agungen liessen mich meine erste Recognoscirungsreise nicht ohne "Angstlichkeit vornehmen (gegen Ende Aprils). Anfangs ging es selbst besser als ich erwartet hatte; ich fand, dass Garssen und Falkenberg sich unmittelbar mit Deister und Lichtenberg verbinden liessen; aber bei den Untersuchungen, wie von jenen beiden Punkten die Dreiecke weiter n"ordlich bis L"uneburg gef"uhrt werden k"onnten, fand ich das Terrain so widerspenstig, dass ich mehrere Male die M"oglichkeit eines gl"ucklichen Erfolgs bezweifelte und bef"urchtete, das ganze Unternehmen aufgeben zu m"ussen. Das Land "uberall flach, keine dominirende Punkte, "uberall Holz, theils in grossen Waldungen, wie der Hassel, der L"using, das Becklinger Holz, etc., theils in unz"ahlbaren kleinern Kampen, die sich schachbrettartig vor einander schieben.

\newpage
\noindent Die Versuche meines Gehilfen, auf der Westseite etwas brauchbares aufzufinden, waren ganz ohne Erfolg; mir selbst gelang es endlich nach den beschwerlichsten Versuchen, gleichsam im Herzen der Heide zwei Dreiecke 9.~10.~12 [Falkenberg-Hauselberg-Wulfsode] und 10.~12.~13 [Hauselberg-Wulfsode-Wilsede] zu etabliren; allein ich "uberzeugte mich zugleich, dass ich bei dieser Gattung von Arbeiten bald unterliegen w"urde, und setzte daher die weitere Aufsuchung einer M"oglichkeit, diese Dreiecke mit den s"udlichern zu verkn"upfen und weiter n"ordlich fortzuf"uhren, auf die sp"atere Zeit hinaus, wo ich st"arkeres Gehilfenpersonal (ich hatte nur einen Officier zur ersten Reise mitgenommen) und alle meine Instrumente bei mir haben w"urde. Ich fing daher Mitte Junius die wirklichen Messungen in Lichtenberg an, und meine Hoffnung ist sp"ater auch so ziemlich erf"ullt. Da es ganz unm"oglich befunden wurde, Garssen vermittelst Durchhaus mit Hauselberg zu verkn"upfen, so wurde nach langem Suchen noch ein n"ordlicherer Punkt 14 [Scharnhorst] gefunden, der mehr Hoffnung darbot.

Allein sp"ater zeigte sich dies absolut unm"oglich, da der Wald zwischen 10 [Hauselberg] und 14 [Scharnhorst] auf zu hohem oder vielmehr nicht genug niedrigem Terrain lag, und ich musste mich gl"ucklich sch"atzen, noch einen andern Punkt, Breithorn, zu finden, womit es endlich gelang. Inzwischen war auch die Linie 9.~13 (Falkenberg-Wilsede) vermittelst eines sehr bedeutenden Durchhaus ge"offnet und so hatte nun 10 [Hauselberg], wenn man 11 [Breithorn] fr"uher gekannt h"atte, ganz wegbleiben k"onnen. Allein ich behielt jenen Punkt mit bei, und halte es f"ur einen "uberaus sch"atzbaren Vortheil, dass in meinem Systeme drei Vierecke vorkommen, in denen alle sechs Richtungen wirklich hin und zur"uck gemessen sind. Nachdem alle Winkel vorher in jedem Dreieck zur geh"origen Summe geh"orig ausgeglichen waren, bedurfte es nur noch sehr kleiner Modificationen, meistens unter $0''\!\!\!,\!51$, um diese drei Vierecke in volle Harmonie zu bringen.

\newpage
\noindent Es w"are zu w"unschen, dass man bei jeder Messung solche Pr"ufungen h"atte. Es gibt Messungen, wobei die Summen der drei Winkel "uberall zum Bewundern stimmen, und wo eine solche Pr"ufung zeigt, dass manche Winkel um $2''$ bis $3''$ gewiss unrichtig sind. In der That ist die Pr"ufung vermittelst der Summe der Winkel \`{a} la port\'{e}e de jedermann; die durch Diagonalen ist es weniger, so leicht sie auch f"ur einen Mathematiker ist, und man kann sich der Vermuthung nicht erwehren, dass die erstere Pr"ufung zuweilen dazu gedient haben mag, wenn auch nicht die Beobachtungen zu verf"alschen, doch etwas zu w"ahlen (man bemerkt eine Tendenz dazu selbst bei Delambre). N"ordlich von Wulfsode ist zum Gl"uck noch der Timpenberg gefunden, der unmittelbar mit Hamburg verbunden werden kann; der versuchte Durchhau von Timpenberg nach L"uneburg ist missgl"uckt, weil das Land dazwischen zu wenig deprimirt war. Die Richtungen von 10 nach 11 und von 11 nach 14 haben grosse Durchhaue erfordert.

Immer machte ich es mir zum Gesetz, mit der Rechnung allen Messungen, wie ich sie erhalten hatte, gleichen Schritt zu halten (bis auf die allerletzte Zeile), und nur dadurch ist es m"oglich geworden, alle Durchhaue mit der "aussersten Pr"acision so durchzuf"uhren, dass auch nicht Ein Stamm ohne Noth gef"allt ist, oder die Unm"oglichkeit der Durchhaue so fr"uh wie m"oglich bestimmt zu erkennen. So wusste ich z.~B.\ die Depression, unter der 14 [Scharnhorst] in 10 (Hauselberg) oder 20 (L"uneburg) in 19 [Timpenberg] erscheinen musste, genau voraus; f"ur den ersten Fall war schon das Terrain vor dem Holz zu hoch, im zweiten, nachdem ein schmaler Spalt 2000 Schritt weit fortgef"uhrt war.

Ich habe von Wilsede aus noch einen Punkt [Nindorf] 3000 Meter n"ordlich von Timpenberg festgelegt, der unmittelbar mit 20 [L"uneburg], 21 [Lauenburg], 23 [Hamburg] communicirt und wahrscheinlich mit 19 [Timpenberg] verkn"upft werden kann, aber die ganze Strecke dahin muss durchgehauen werden. Auf Timpenberg habe ich die Winkel zwischen 12 [Wulfsode], 13 [Wilsede], 23 [Hamburg] erst vorl"aufig gemessen, alle s"udlichern Punkte sind absolvirt sowie Wilsede.

\newpage
\noindent Bei Scharnhorst bin ich zuletzt gewesen; dieser Punkt liesse sich vermittelst leichter Durchhaue mit Lichtenberg und Deister unmittelbar verbinden, wodurch Garssen entbehrlich w"urde; allein es w"are unm"oglich gewesen, die Brauchbarkeit jenes Punkts zu Anfang auszumitteln. Dieser Punkt wie mehrere andere sind ganz unscheinbare Pl"atze, von denen man gar nicht vermuthen sollte, dass sie so vielen Werth haben. Bei vielen Richtungen, z.~B.\ von Lichtenberg nach Garssen und von Lichtenberg nach Falkenberg, geht die Linie so knapp "uber die Zwischenhindernisse, dass (der Zielpunkt) bei gew"ohnlicher Refraction nur wenige Secunden sich dar"uber erhebt und nur zuweilen $30''$ bis $40''$ erreichte; von Falkenberg nach Wulfsode und \emph{vice versa} kam das Licht bei schw"acherer Refraction gar nicht her"uber und hob sich immer erst in den sp"atern Nachmittagsstunden herauf.

Vom Lichtenberg selbst sah ich auf dem Falkenberg selten etwas und \emph{vice versa}, das Heliotroplicht schwebte fast immer im freien Himmel (Distanz 85542 Meter); der grossen Distanz ungeachtet stimmen die Messungen ebenso gut oder vielmehr fast besser als auf ganz kleinen Distanzen von 10000 Meter Entfernung, wo das Heliotroplicht noch immer fast zu stark ist, wenngleich der Spiegel bis auf eine "Offnung von wenigen Quadratmillimetern bedeckt war. Um an Zeit zu gewinnen, habe ich "Ofters auch selbst in Distanzen von 3 bis 4 bis 5 Meilen auf meine steinernen Postamente, $3\frac{1}{2}$ Fuss hoch, selbst pointirt. Dadurch sind die Messungen hin und wieder etwas weniger genau geworden, als wenn ich bloss Heliotroplicht gebraucht h"atte, allein dann w"are ich in diesem Jahre lange nicht so weit gekommen. Der gr"osste Fehler der Summe der drei Winkel war in diesem Jahre $1\frac{5}{7}''\!,\!6$, auch "uberhaupt\footnote{In 19 Dreiecken.} der gr"osste nachst dem

\newpage
\noindent im Dreieck 4.~5.~6 [Hohehagen-Hils-Brocken], wo er $3,\!7$ betr"agt und haupts"achlich dem schwierigen Pointiren auf den Brockenhaus-Thurm zuzuschreiben ist, welcher bei Sonnenschein nie gut geschnitten werden kann wegen der Phase. Gern m"asse ich die Winkel dieses Dreiecks, wenn es die Zeit erlaubte, noch einmal nach.

In diesem Jahre hatte ich oft drei Heliotrope zugleich in Activit"at, deren einen mein Sohn besorgte; es ist in der That ein prachtvoller Anblick. Die Nachricht, welche ein gewisser Schupach "uber die Heliotrope im [Astronomischen] Jahrbuch 1825 gegeben hat, beruht auf einem Irrthum und hat gar nichts mit meinen Heliotropen gemein; diese sind k"unstliche Instrumente, deren Einrichtung mir erst sehr viele M"uhe gekostet hat, die nun aber auch, wie ich glaube, nichts zu w"unschen "ubrig lassen. Ich habe von beiden Einrichtungen, die unter sich ganz verschieden sind, Exemplare f"ur den General M"uffling hier anfertigen lassen, auch Gerling hat zwei erhalten; vielleicht kann ich bald in Schumachers Astronomischen Nachrichten Abbildungen davon geben. Zum Telegraphiren und um bei grossen Entfernungen den gegen"uberstehenden Heliotropen erst die Richtung zu zeigen, brauchte ich oft einen grossen Spiegel von einem Fuss Quadrat, welcher wieder auf andere Art gelenkt wurde. Der Anblick davon ist nach der Beschreibung meiner Gehilfen h"ochst prachtvoll gewesen; auf vier Meilen weit hat es dem blossen Auge zuweilen wehe gethan, lange hinzusehen. --- Doch jetzt genug hievon.

Ihre Art, geod"atische Beobachtungen zu behandeln, habe ich mit Vergn"ugen in Schumachers Astronomischen Nachrichten gesehen. Sie wissen, dass dieser Gegenstand mich schon vor vielen Jahren besch"aftigt hat. Da Ihren Arbeiten nicht leicht etwas beigef"ugt werden kann, so w"urde, h"atten unsere Wege sich begegnet, jene Bekanntmachung meine eigene Arbeit "uberfl"ussig gemacht haben. Allein die Art, wie ich diesen Gegenstand behandelt habe, ist von der Ihrigen durch und durch verschieden, und so werde ich also in Zukunft bei Bekanntmachung meiner Messungen auch meine theoretischen Arbeiten ausf"uhrlich entwickeln. Ich hoffe darin manches unerwartete geben zu k"onnen. Aber diese Untersuchungen hingen mit einem reichen, fast unersch"opflich reichen Felde zusammen, und ich f"uhle oft mit inniger Wehmuth, bei dieser wie bei so vielen andern Gelegenheiten, wie meine "aussern Verh"altnisse mich an weitaussehenden theoretischen Arbeiten hindern. Wenn solche ganz gedeihen sollen, muss man sich ihnen ganz hingeben k"onnen und nicht durch so heterogene Arbeiten wie Collegia lesen, alles kleinliche Detail beim Observiren und Rechnen der Beobachtungen, etc.\ etc.\ st"undlich gehindert werden.\,\ldots\ldots\ldots

\newpage
\noindent Die F"atigen im heissen Sommer sind oft "ausserst angreifend f"ur mich gewesen, zuweilen so, dass ich glaubte, ich w"urde ihnen erliegen. Auch das ist eine grosse Beschwerde bei den Arbeiten in der "oden L"uneburger Heide, dass man "ofters nur ein schlechtes Unterkommen und doch selbst ein solches nur meilenweit vom Arbeitspunkte haben kann. Bei k"uhlem Wetter, welches meiner Constitution besser zusagt, befand ich mich im allgemeinen immer leidlich wohl, und jetzt kann ich "uber mein Befinden nicht klagen.\,\ldots\ldots\ldots

Sie k"onnen aus obigem Bericht selbst sehen, was zur Vollendung meiner Triangulirung noch fehlt, aber einen eigentlichen Plan f"ur n"achsten Sommer kann ich jetzt noch nicht machen, es wird dabei auch vieles auf Schumachers Cooperation ankommen. Es w"are m"oglich, dass ich, wenn ich mich bloss auf das Indispensable beschr"anke, schon im April und Mai die Dreiecke beendige und dann hieher zur"uckkomme; vielleicht im Julius hier und im August an einem n"ordlichern Punkte (Celle, Harburg, Hamburg?) mit dem Zenithsector messe. Es k"onnte aber auch sein, wenn ich die Kosten nicht zu scheuen brauche, dass ich noch den ganzen Sommer auf die Triangulirung wende, um alles zur m"oglich gr"ossten Vollkommenheit zu bringen.

Vor einigen Tagen habe ich aus M"unchen ein Universalinstrument erhalten, das ich vor zwei Jahren vorz"uglich behuf der Azimuthe bestellt hatte. Inzwischen scheint nach meiner vorl"aufigen Reduction das von mein Passageninstrument entlehnte und bis Hamburg "ubertragene Azimuth von dem, was mir Schumacher mitgetheilt hat (er hat auch vorl"aufig den Winkel zwischen L"uneburg und Wilsede, von wo ich ihm Heliotroplicht zusenden liess, gemessen), noch nicht $1\frac{3}{5}''$ zu differiren und also eine Azimuthmessung von meiner Seite an andern Punkten wohl ziemlich "uberfl"ussig zu sein. Ob ich mich veranlasst sehen werde, mit dem Universalinstrument auch Polh"ohen an mehrern Punkten zu messen, wird sich zeigen, wenn ich das Instrument erst eine Zeit lang gebraucht habe. Es ist gut gearbeitet, scheint mir aber f"ur den Gebrauch nicht recht bequem zu sein, besonders f"ur ein kurzsichtiges Auge, welches zum Ablesen bei den Stellkreisen nicht gut zu kann. Auf alle F"alle muss es sehr erm"udend sein, viel damit zu beobachten. Meine Horizontalwinkel habe ich alle mit einem 12-z"olligen Theodolithen gemessen, der ein sehr vortreffliches Instrument ist. Mit dem Zenithsector habe ich im vorigen Winter nur wenige Messungen gleichsam zur Probe gemacht; optische Kraft und Bequemlichkeit des Gebrauchs stehen dem Reichenbachschen Meridiankreis sehr nach.

\newpage
\noindent Meine Gradmessung als solche kann eigentlich f"ur sich allein kein sehr wichtiges Resultat liefern; ich weiss aber nicht, wann die Vollendung der Schumacherschen zu hoffen ist. So viel ich weiss, hat er seine Dreiecke erst etwa ein Drittel der ganzen L"ange gef"uhrt, auch die Zenithsector-Messungen in Skagen sind bloss von seinem Gehilfen Caucho gemacht. Die genaue L"ange seiner Basis (circa 6000 Meter) kenne ich noch nicht; erst in diesem Herbst wollte er sie mit Hamburg verbinden, in seinem letzten Briefe erw"ahnt er aber gar nichts davon. Ich k"onnte nun zwar selbst eine Basis messen und die Linie von Breithorn bis Scharnhorst scheint in ihrer ganzen L"ange (11220 Meter) keine un"ubersteigliche Hindernisse darzubieten. Allein auch abgesehen von den grossen Kosten gestehe ich, mich vor einer so h"ochst langweiligen Arbeit zu scheuen. Meinen trigonometrischen Messungen habe ich immer eine interessante Seite abgewinnen k"onnen, da ihre t"agliche Reduction immer einige Unterhaltung gab.\footnote{Meine Beobachtungsmanier war, immer zu messen, was sich eben gut messen liess, ohne R"ucksicht, ob es ein unmittelbarer Dreieckswinkel war. An manchen Stationen nahm ich einige Halfspunkte, um auch dann nicht missig zu sein, wenn nur Ein Heliotrop leuchtete. Die Franzosen haben, wie ich sehe, etwas "ahnliches in Spanien gethan, aber die Messungen ganz unrichtig ausgeglichen.} Ich schnitt "uberdies auch alle sichtbaren Objecte bei Gelegenheit, um mich f"ur die Landesgeographie n"utzlich zu beweisen, und ich muss sagen, dass ich dieses Gesch"aft mit seinen t"aglichen Ausgleichungen so lieb gewann, dass das Bemerken, Ausmitteln und Berechnen eines neuen Kirchthurms wohl ebenso viel Vergn"ugen machte, wie das Beobachten eines neuen Gestirns. (Vor Gott ist's am Ende auch wohl einerlei, ob wir die Lage eines Kirchthurms auf einen Fuss oder die eines Sterns auf eine Secunde bestimmt haben.) Allein bei einer Basismessung, sobald sie einmal im Gange ist, sehe ich f"ur den Verstand auch gar nichts, was ihn reizen oder unterhalten k"onnte, und man muss sich bei einer vielleicht zwei Monate dauernden angestrengten t"aglichen Arbeit lediglich mit dem Gedanken aufrecht halten, dass es eben doch einmal geschehen muss, um zuletzt eine Zahl zu haben.

\newpage
\noindent Die grosse Genauigkeit im Messen horizontaler Winkel durch Heliotroplicht auf die ungeheuersten Distanzen, und die Genauigkeit, womit man durch Universalinstrumente absolute Azimuthe messen kann, lassen mich glauben, dass man gegenw"artig eine L"angengradmessung in schicklichem Terrain mit viel Vortheil ausf"uhren k"onnte, indem man den L"angenunterschied nicht auf Zeitbestimmung, sondern auf die Convergenz der Meridiane gr"undete. Von den Pyren"aen sieht man den Mont Ventoux, von diesem die Alpen, von diesen bis Steiermark, von da bis Ungarn, etc. Freilich w"achst hiebei die Genauigkeit nicht wie beim Meridianbogen, wie die ganze L"ange des Bogens (den eigentlich geod"atischen Theil kann man dabei als fehlerfrei betrachten); also wenn man sich St"ucke von ungef"ahr gleicher Gr"osse denkt, w"achst beim Breitengrade die Genauigkeit wie die Anzahl der St"ucke, beim L"angengrade ebenfalls, wenn der L"angenunterschied der Endpunkte astronomisch, z.~B.\ durch Sternbedeckungen, beobachtet wird; hingegen bei der Convergenz der Meridiane nur wie die Quadratwurzel der Zahl der St"ucke (ungef"ahr); allein ich glaube doch, dass bei sehr grossen St"ucken dies Verfahren mehr Genauigkeit gibt, als man z.~B.\ erhalten kann, wenn man den L"angenunterschied durch Pulversignale vermittelst Zeitbestimmung sucht. Am vortheilhaftesten w"are dies Verfahren in n"ordlichen Gegenden, wenn es dort sehr grosse Fernsichten gibt. Sollten diese nicht in Norwegen und Schweden oder in einigen Gegenden von Russland zu finden sein?

\medskip
\noindent Ich bin jetzt noch beim Ausgleichen, was ich leider einmal ganz umsonst gemacht habe, da durch ein Versehen von einer Station ein fehlerhaftes Tableau aufgenommen war. Bei meiner Behandlung reagirt gewissermassen jeder z.~B.\ in Wilsede gemessene Winkel auf alle "ubrigen bis G"ottingen hin. Ich werde Ihnen k"unftig einige Resultate anzeigen. Nach vorl"aufiger Rechnung, G"ottingen zu $51^{\circ}31'48\frac{7}{10}''$ angenommen, f"allt Hamburg in $53^{\circ}33'15\frac{6}{10}''$ und $0^{\circ}\,2'25\frac{97}{100}''$ "ostlich von G"ottingen, das Absolute vorl"aufig auf Zach's Basis gest"utzt.

\newpage
\noindent \textbf{Gauss an Bessel.} G"ottingen, 5.~November 1823.

\medskip
\noindent Ich habe einen Theil des Jahrs damit zugebracht, den noch "ubrigen Theil meiner trigonometrischen Messungen im Norden: von Timpenberg, Nindorf, L"uneburg bis Hamburg zu absolviren; dann auch vorl"aufig die weiter westlich liegende Gegend, nach Bremen zu, zu recognosciren, da unser Gouvernement eine weitere Fortsetzung der Messungen nach Westen zu, bis zur holl"andischen Grenze und zum Anschluss an die Krayenhoffschen Dreiecke, w"unscht; endlich zuletzt habe ich noch einmal den Brocken und Hohehagen besucht, da theils die Winkel des Dreiecks Hohehagen-Hils-Brocken im Jahr 1821 unter sehr ung"unstigen Umst"anden gemessen, theils jetzt noch die damals missgl"uckte Verbindung des Brockens mit dem Inselsberg zu effectuiren war, sowie auch jetzt noch der hessische Dreieckspunkt Meisner angekn"upft werden sollte. Ich habe diese Zwecke meistens zu meiner Zufriedenheit erreicht; nur den Hils hatte ich gern auch noch einmal besucht, um den Winkel dort genauer zu messen; die vorger"uckte Jahreszeit hat mich aber daran gehindert.

Diesen Umstand abgerechnet, kann ich jetzt die Triangulirung zur Gradmessung, so weit sie zu meinem Ressort geh"ort, als geendigt ansehen. Astronomische Beobachtungen sind noch keine weiter gemacht, als die zur Orientirung meiner ersten Dreiecksseite geh"oren. Ob ich den oben erw"ahnten Plan der Fortsetzung der Messungen nach Westen noch ausf"uhre, ist "ubrigens noch sehr ungewiss. Es ist manches daf"ur, manches, fast noch mehr, dagegen, auch abgesehen davon, dass vielleicht noch die M"oglichkeit einer "Anderung meiner "aussern Lage eintreten k"onnte\footnote{Es handelte sich um die Berufung nach Berlin.}.

\medskip
\noindent Ich habe das System meiner Hauptdreiecke in diesen Tagen sorgfaltig ausgeglichen, so dass nicht nur die Summe der Winkel jedes einzelnen Dreiecks, sondern auch die Verh"altnisse der Seiten in den gekreuzten Vierecken und F"unfecken genau harmoniren, und zwar ohne alle Willk"ur, ohne Ausw"ahlen, ohne Ausschliessen, alles nach der Strenge der Probabilit"atsrechnung. Es sind zusammen 26 Dreiecke, worin alle Winkel von mir selbst beobachtet sind.

\newpage
\noindent Die gr"osste Summe der Fehler ist $2\frac{5}{2}''$ in einem Dreiecke, wo bei einer Seite das Pointiren sehr schwierig war; die n"achst gr"osste ist $1\frac{3}{8}''$. Keine der 76 vorkommenden Richtungen ist bei der Ausgleichung um eine ganze Secunde ge"andert; die gr"osste "Anderung betr"agt $0\frac{7}{8}\!\!\!,\!13$ bei der oben erw"ahnten Seite von Nindorf nach Hamburg. Was ich nach meiner neuen Probabilit"atstheorie den mittlern Fehler nenne, bei den Richtungen, ist $0\frac{5}{4}\!\!\!,\!8$.

\medskip
\noindent \textbf{Gauss an Bessel.} G"ottingen, 20.~November 1824.

\medskip
\noindent Ich habe in diesem Jahre 12 Stationen besucht und bin mit den Dreiecken bis an die Weser (der entfernteste Punkt auf der Garlster Heide zwischen Osterholz und Vegesack) gekommen. Die dritten Winkelpunkte liess ich anfangs immer zur"uck, und als ich nach der Mitte Augusts Bremen verliess, hatte ich noch 6 Pl"atze zu besuchen. Mit allen ging es noch ertr"aglich, aber bei dem letzten, dem Wilseder Berge, qu"alte mich das schlechte Wetter so, dass ich trotz einem dreiw"ochentlichen Aufenthalt nicht ganz zu meiner Zufriedenheit fertig wurde; ich musste zuletzt einen Schlusstermin setzen und kam Ende Octobers nach G"ottingen zur"uck.

Ich behalte mir vor, Ihnen k"unftig ausf"uhrlicher "uber den Erfolg der Messungen, die grossen Schwierigkeiten, mit denen ich zu k"ampfen gehabt (eine davon, das ewige Moorbrennen, wird Ihnen selbst noch im Ged"achtniss sein), und manche interessante Ph"anomene, die sich dabei ergeben haben, zu schreiben.\,\ldots\ldots\ldots

\medskip
\noindent \textbf{Gauss an Bessel.} G"ottingen, 12.~M"arz 1826.

\medskip
\noindent Den gr"ossten Theil des vorigen Sommers habe ich im Bremischen und Oldenburgschen mit meinen Messungen zugebracht.\,\ldots\ldots\ldots

Was meine Messungen betrifft, so habe ich deren trigonometrischen Theil, wenigstens dem Buchstaben nach, vollendet; meine Seite Varel-Jever schliesst sich an die Krayenhoffschen Dreiecke. Ob sie aber wirklich geendigt sind, weiss ich selbst noch nicht. Meine Winkelmessungen in Jever geben ganz enorme Unterschiede von den Krayenhoffschen, die bis auf $15''$ gehen. Ihre Quelle kann ich nicht mit Gewissheit angeben. Meine eigenen Winkel verb"urge ich bis auf eine Secunde. Die Differenzen w"urden sich erkl"aren lassen, wenn ich annehmen d"urfte, dass Krayenhoff's Centrum der Station von dem meinigen (ich habe eine Etage tiefer beobachtet) ein Meter entfernt liegt. Die Aufschl"usse, die ich von Krayenhoff erhalten habe, sind unbefriedigend in R"ucksicht auf das Centriren; im allgemeinen aber geht daraus hervor, dass seine Winkelmessungen in dortiger Gegend lange nicht die Sch"arfe haben wie die meinigen. Er hat in Ostfriesland ein schlechteres Instrument gebraucht; aus vielen Winkelreihen hat er immer nur diejenigen beibehalten, die am besten zu passen schienen (ohne anzugeben, wie viel die andern abwichen) und selbst unter den beibehaltenen in Jever finden sich Differenzen von $4\frac{7}{7}''$. Unter diesen Umst"anden scheint es mir nicht rathsam, meine Messungen in Ostfriesland weiter auszudehnen. Seine s"udlichern Messungen sind besser, und eine Verbindung meiner Dreiecke "uber das Osnabr"ucksche nach Bentheim w"urde ohne Zweifel zuverl"assigere Resultate geben k"onnen. Allein der wankende Zustand meiner Gesundheit hat mich bisher muthlos gemacht, auf solche Operationen, die noch eine ein- oder anderthalbj"ahrige Campagne erfordern w"urden, anzutragen, und in diesem Sommer wird schwerlich etwas erhebliches darin geschehen k"onnen.\,\ldots\ldots\ldots

\newpage
\noindent \textbf{Gauss an Bessel.} G"ottingen, 20.~November 1826.

\medskip
\noindent Die Verarbeitung der Materialien zu dem beabsichtigten Werke "uber meine Messungen kostet mich viele Zeit. Meine Hauptdreiecke, 33 Punkte befassend, sind zwar l"angst fertig berechnet, aber die Berechnung der vielen geschnittenen Nebenpunkte, die f"ur die Geographie eines bedeutenden Theils von Norddeutschland wichtig sind, macht viel Arbeit, da jene bisher entweder noch gar nicht oder nur provisorisch berechnet waren. Mein Verzeichniss enth"alt jetzt etwa 250 Punkte, alle aus dem n"ordlichern Theil meiner Messungen. Da meine Dreiecke durch die kurhessischen mit den bayerischen und w"urttembergischen zusammenh"angen, und letztere mir auch von Bohnenberger gef"allig mitgetheilt sind (ebenso wie die darmst"adtschen von Eckhart), so ist es mir sehr unangenehm, dass alle meine Bem"uhungen, die bayerischen zu erhalten, bisher vergeblich gewesen sind.

Schumacher erz"ahlte mir, dass Soldner ihm gesagt h"atte, der Grund, warum die Commission in M"unchen meine Bitte um die Mittheilung nicht erf"ullt habe, sei, weil man annehme, dass ich "uber diese Dreiecke Rechnungen anstellen wolle! Ebenso ist mir's mit den "osterreichischen gegangen. Auch General von M"uffling sollte doch seine Dreiecke "ostlich von Seeberg, die Berlin anschliessen und sich bis "uber Schlesien erstrecken, bekannt machen.

Noch viel mehr Verlegenheit macht mir der weit ausgedehntere theoretische Theil, der so vielfach in andere Theile der Mathematik eingreift. Ich sehe hier kein anderes Mittel, als mehrere grosse Hauptpartheien von dem Werke abzutrennen, damit sie selbstst"andig und in geh"origer Ausf"uhrlichkeit entwickelt werden k"onnen. Gewissermassen habe ich damit schon in meiner Schrift "uber die Abbildung der Fl"achen unter Erhaltung der "Ahnlichkeit der kleinsten Theile den Anfang gemacht; eine zweite Abhandlung, die ich vor ein paar Monaten der k"oniglichen Societ"at "ubergeben habe, und die hoffentlich bald gedruckt werden wird, enth"alt die Grunds"atze und Methoden zur Ausgleichung der Messungen (beil"aufig ist daraus indirect auch ersichtlich, wie weit die Krayenhoffschen Messungen von derjenigen Genauigkeit entfernt sind, die man ihnen mit Unrecht beigelegt hat). Vielleicht werde ich zun"achst erst noch eine dritte Abhandlung ausarbeiten, die mancherlei neue Lehrs"atze "uber krumme Fl"achen, k"urzeste Linien, Darstellung krummer Fl"achen in der Ebene, u.\ s.\ w.\ entwickeln wird. H"atten alle diese Gegenst"ande in mein projectirtes Werk aufgenommen werden sollen, so h"atte ich entweder manches ungr"undlich abfertigen oder dem Werk ein sehr buntscheckiges Ansehen geben m"ussen.\,\ldots\ldots\ldots

\newpage
\noindent \textbf{Gauss an Bessel.} G"ottingen, 1.~April 1827.

\medskip
\noindent Ich denke in diesem Fr"uhjahr die Amplitudo des Bogens zwischen G"ottingen und Altona mit dem Zenithsector zu messen und bin selbst begierig auf das Resultat. Eine Reihe Beobachtungen, die an den Meridiankreisen Anfangs 1824 gemacht war, gab $2^{\circ}1'58''$; die Polh"ohen, wie ich die meinige nach den besten Beobachtungen annehme und wie Schumacher die seinige angibt, geben eher 1 oder $1\frac{1}{2}$ Secunden weniger; die geod"atischen Messungen hingegen unter Voraussetzung der gleichf"ormigen Gestalt der Erde und Walbeck's Dimensionen geben 4 oder 5 Secunden mehr, und fast genau ebenso viel, wie dies letztere Resultat, geben ganz gleichzeitige Zenithdistanzen von Zenithalsternen, die ich hier und Repsold in Altona am Meridiankreise beobachtet haben, indem der Nullpunkt mit Collimatoren bestimmt war. Repsold ist jetzt hier, mir bei den Sectorbeobachtungen zu helfen; er hat einen Repsold'schen pensilen Collimator mitgebracht, bis jetzt aber halte ich die Methode des Nadirpunkts durch Quecksilberreflexion f"ur bedeutend genauer; doch, denke ich, wird jene Methode sich noch verbessern lassen.\,\ldots\ldots\ldots

\medskip
\noindent \textbf{Gauss an Bessel.} G"ottingen, 9.~April 1830.

\medskip
\noindent Noch viel mehr Zeit haben mir seit Mai 1829 die trigonometrischen Messungen geraubt, wenn ich gleich keinen unmittelbaren Antheil an den Gesch"aften im Felde diesmal genommen habe. Noch diese Stunde bin ich nicht ganz (obwohl Gott Lob beinahe) mit Verarbeitung der vorigj"ahrigen Messungen fertig, wobei ich jeder H"ulfe entbehre. Doch habe ich dabei viel Freude "uber die zu meiner gr"ossten Zufriedenheit ausgefallene Art gehabt, wie mein Sohn seinen Antheil an diesen Gesch"aften ausgef"uhrt hat. Er hat ganz allein ein grosses Dreiecksnetz von der Weser bis zur holl"andischen Grenze ("uber das Osnabr"ucksche, auch mit mehrmaliger Ber"uhrung preussischen Gebiets, wobei ihm von Seiten der dortigen Beh"orden sehr liberal Vorschub geleistet ist) gef"uhrt, wodurch ausser den Hauptpunkten noch gegen 250 andere festgelegt sind. Die --- von einem andern Officier --- im F"urstenthum Hildesheim gemachten trigonometrischen Messungen sind gleichfalls so gut wie vollendet, so dass alle Messtischbl"atter (53, wovon etwa ein Viertel bereits aufgenommen ist) mit festen Punkten versehen werden k"onnen.\,\ldots\ldots\ldots

\newpage
\subsection*{[3.]}
\addcontentsline{toc}{subsection}{Gauss an Bohnenberger}

\noindent \textbf{Gauss an Bohnenberger.} G"ottingen, 16.~November 1823.

\medskip
\noindent Mit Vergn"ugen habe ich aus Ihrem Briefe einiges von Ihren Messungen erfahren, und m"ochte nichts lieber, als dass diese Antwort Veranlassung geben m"ochte, etwas ausf"uhrlicheres dar"uber mitgetheilt zu erhalten.

Was zuerst die Heliotrope betrifft, so habe ich zwei ganz verschiedene Arten anfertigen lassen, die auf ganz verschiedenen Principien beruhen; die zweite Einrichtung finde ich aber am vortheilhaftesten.\,\ldots\ldots\ldots\ Durch die M"oglichkeit, wo es sonst das Terrain erlaubt, die gr"ossten Dreiecke anzuwenden, "uberall gleich anfangen zu k"onnen, ohne erst die so viel Zeit und Geld kostenden Signale errichten zu m"ussen, wird die kleine Ausgabe [f"ur die Heliotrope] vielfach erspart, obwohl dies der geringste Vortheil ist; die Messungen werden dadurch einer Sch"arfe f"ahig, auf die man bei Signalen und Kirchth"urmen selten rechnen kann. Meine schlechtesten Dreiecke (relativ gesprochen) sind die, worin Th"urme die Zielpunkte waren. So viel von den Heliotropen.

Was die Messungen selbst betrifft, so w"unsche ich nichts sehnlicher, als bald die Verbindung mit Ihnen zu haben. Ich habe jetzt meine Dreiecke mit den kurhessischen zusammengeh"angt, wovon beiliegende Zeichnung Ihnen einen Begriff gibt\footnote{Eine Zeichnung liegt der im Gauss-Archiv befindlichen Copie des Originals nicht bei.}. Zwar ist von den hessischen Dreiecken nur erst ein Theil wirklich gemessen, aber schon genug, um alle Punkte bis zum Feldberg mit den meinigen wenigstens vorl"aufig zu verkn"upfen. Meine n"ordlicher liegenden Dreiecke kennen Sie aus Schumachers Astronomischen Nachrichten, I.\ Nr.~24, die in diesem Jahr noch dazu gekommenen n"ordlichen Punkte werden Sie wenigstens f"ur den Augenblick nicht interessiren.

Ich vermuthe nun, dass hiedurch und die darmst"adtschen und die bayerischen Dreiecke meine Messungen mit den Ihrigen verbunden sind, besitze aber von den letztern noch gar keine, von den erstern nur einige fragmentarische Angaben. Sollten Sie in vollst"andigerm Besitze sein, so verpflichten Sie mich ausserordentlich durch Mittheilung, ebenso wie von Ihren eigenen Dreiecken. Nur w"unschte ich die Winkel zwar auf die Centra reducirt, aber ohne Fehlerausgleichung zu erhalten\footnote{"Ubrigens allenfalls ohne alle weitern Rechnungsresultate.}. Ich m"ochte gern alles nach gleichf"ormiger Methode berechnen. Die Verkn"upfung mit T"ubingen, Mannheim und M"unchen interessirt mich um so dringender, da, im Vertrauen gesagt, die ersten Messungen Schumachers in Altona, wo er eine sch"one kleine Sternwarte mit einem dem meinigen, Soldner'schen und Bessel'schen ganz gleichen Reichenbachschen Meridiankreise errichtet hat, eine Polh"ohe gegeben haben, die $5''$ von der von G"ottingen durch die Dreiecke "ubertragenen differirt. Bei der Rechnung habe ich die Dimensionen der Erde, die Walbeck aus dem Ensemble aller guten Gradmessungen abgeleitet hat, gebraucht:
\begin{align*}
a &= 3272077.6 \quad \text{[Toisen]} \\
b &= 3261467.2 \quad \text{[Toisen]}
\end{align*}
Abpl.\ $= \frac{1}{302.78}$\,\footnote{Siehe die Briefe an Olbers vom 6.~Julius 1824, S.~320 Anmerkung, und vom 1.~M"arz 1827, S.~378.}. Uebrigens aber sind meine Rechnungsmethoden so g"anzlich von den sonst angewandten verschieden, dass ich in einem Briefe Ihnen keinen Begriff davon geben kann. Ich habe die Absicht, wenn meine Messungen erst vollendet sind, diese Methoden zu einem gr"ossern Werke zu verarbeiten und durch Anwendung auf die hannoverschen und die damit zusammenh"angenden Messungen zu erkl"aren. Unser Gouvernement ist geneigt, die hannoverschen Messungen weiter westlich auszudehnen, wodurch sie mit den Krayenhoffschen und dadurch mit den franz"osischen und englischen in Zusammenhang kommen. Krayenhoff's Messungen sind bekanntlich gedruckt; dies sollte mit allen guten Triangulirungen geschehen; die grossen Dreiecke geh"oren gewissermassen der ganzen cultivirten Mit- und Nachwelt an, um so mehr, je mehr sie nach und nach unter sich in Zusammenhang kommen.

\newpage
\noindent Die M"ufflingschen Dreiecke h"angen zwar mit der franz"osischen Gradmessung auch schon durch Tranchot's Dreiecke zusammen; allein diese sind nicht gedruckt, also f"ur das Publicum so gut wie gar nicht vorhanden, und General von M"uffling selbst besitzt sie nur in ungen"ugender Form (vermuthlich die fatalen Chordenwinkel zu $180^{\circ}$ schon abgeglichen). Leider finde ich, dass die Menschen so wenig zur Communication geneigt sind; ich habe mir auf officielem diplomatischen und auf nicht officielem Privatwege viele M"uhe gegeben, aus Paris die von Enky 1804 und 1805 im Hannoverschen gemachten Messungen zu erhalten, aber nichts als Ausfl"uchte, eine blosse Namenangabe der Stationen und eine Zeichnung der Dreiecke erhalten, 2 oder 3 Zahlangaben nicht gerechnet, die, wie aus meinen Messungen folgt, entschieden grob unrichtig sind. Lassen Sie uns eine Ausnahme davon machen. Ich wiederhole nochmals meine Bitte um eine Mittheilung Ihrer Messungen, und um freundschaftliche Mittheilung dessen, was Sie von fremden mit den unsrigen zusammenh"angenden haben, insofern ich es direct nicht erhalten kann, und erbiete mich gern \emph{ad reciproca}.

Wie sch"on w"are es, wenn einmal alle "uber Europa, von Schottland bis zum Banat und von Copenhagen bis Genua und Formentera, sich erstreckenden Messungen in Ein zusammenh"angendes System gebracht werden k"onnten. Ich m"ochte gern nach Kr"aften dazu vorbereiten, allein wenn man "uber die Mitte seines Lebens hinaus ist, muss man bei einem so ausgedehnten Gegenstand je eher je lieber anfangen.

Gerling's Messungen sind mit einem 12-z"olligen Ertel'schen Theodolithen gemacht, ganz dem meinigen und Schumacher'schen gleich\footnote{Derselbe hat auch 3 Heliotrope, von Herrn Rumpf, nach der zweiten Einrichtung.}. Bei meinen Messungen habe ich gefunden, dass das, was ich in meiner Abhandlung in den neuesten G"ottinger Commentationes \emph{>>Theoria combinationis observationum etc.<<} den \emph{mittlern Fehler} nenne, aus mehrern Stationen, gute und weniger gute Messungen durch einander gerechnet, etwa $= \frac{a}{n}$ ist, $n =$ Anzahl der Repetitionen. Bei sehr fester Aufstellung, sehr g"unstiger (d.\ i.\ nicht zitternder) Luft und ausschliesslich heliotropischen Zielpunkten ist er aber betr"achtlich kleiner.

Meine s"ammtlichen Messungen geben bisher 76 Hauptrichtungen (38 hin und 38 zur"uck) und aus der Ausgleichung der Fehler fand sich, dass der mittlere Fehler einer Hauptrichtung $= 0\frac{5}{4}\!\!\!,\!7$ war. Es bilden sich daraus zusammen 26 Dreiecke, worin alle Winkel von mir gemessen sind; darunter mehrere, die gekreuzte Vierecke und F"unfecke geben, und die Ausgleichung ist ohne Willk"ur, ohne Ausw"ahlen und ohne Ausschliessen gemacht, nach strengen Gr"unden der Wahrscheinlichkeits-Rechnung, so dass zuletzt alles genau zu einander passt. Solche gekreuzte Vierecke w"urden bei manchen Messungen ein trefflicher Probirstein sein, wo man findet, dass die Summen der Winkel zwar "uberall vortrefflich passen, so dass selten ein Dreieck viel "uber $1''$ fehlt, wo aber jene Pr"ufung (die nicht in dem Grade \`{a} la port\'{e}e de jedermann ist, wie das Berechnen eines sph"arischen Excesses), wenn sie angewandt werden kann, zuweilen ganz entschieden zeigt, dass Fehler von $2''$, $3''$ oder dar"uber in einzelnen Winkeln vorhanden sind. Der gr"osste Fehler in der Summe der noch nicht ausgeglichenen Winkel bei meinen 26 Dreiecken war $2\frac{5}{2}''$, wo eine Richtung auf den sehr schwer zu schneidenden und bei Sonnenschein nicht ganz phasenfreien Michaelisthurm in Hamburg ging; der n"achst gr"osste $1\frac{7}{8}''$ in einem Dreiecke, wo auch eine Richtung auf einen "ausserst schwer zu sehenden nicht heliotropischen Zielpunkt ging. Ich h"atte gew"unscht, die letztere betreffende Station noch einmal zu besuchen und die Richtung durch Heliotroplicht zu nehmen, konnte aber in diesem Jahre nicht mehr dazu kommen.

Das grosse Dreieck Hohehagen-Brocken-Inselsberg ist unter den 26 nicht begriffen, der Winkel auf dem Inselsberg ist von Gerling gemessen, und meinen Messungen auf dem Brocken war das Wetter sehr ung"unstig, so dass ich den Inselsberg nur 15-mal habe schneiden k"onnen (1823); verbunden mit den 15 Schnitten von 1821, die weniger als $1''$ von jenen differiren, geben sie aber doch einen vortrefflichen Schluss dieses grossen Dreiecks. --- In einem Lande wie W"urttemberg und Kurhessen, wo es so viele hohe Punkte gibt, ist das Messen ein Vergn"ugen, und die gr"ossten Dreiecke leicht aufzufinden. Ebenso im s"udlichen Theile des Hannoverschen. Aber im n"ordlichen, der L"uneburger Heide, habe ich uns"agliche Schwierigkeiten gehabt, und eine im vorigen Sommer nach Westen, gegen Bremen zu, unternommene Recognoscirung hat noch keine Resultate f"ur die M"oglichkeit nur leidlich guter Dreiecke gegeben. Durch hohe Ger"uste liessen sich die Schwierigkeiten zwar wohl "uberwinden, aber ich f"urchte mich vor den grossen Kosten an Geld und Zeit und noch mehr vor der Einbusse solider Aufstellung. Weiter s"udlich durch das Osnabr"ucksche nach Bentheim liesse sich vermuthlich leichter durchkommen, die Messungen w"urden dann aber dem gr"ossten Theile nach "uber fremdes Gebiet gehen m"ussen.

\newpage
\subsection*{[4.]}
\addcontentsline{toc}{subsection}{Gauss an Olbers}

\noindent \textbf{Gauss an Olbers.} G"ottingen, 13.~Januar 1821.

\medskip
\noindent Sollte es in Zukunft bei der wirklichen Messung mit allen "aussern Umst"anden besser gehen, als es bisher den Anschein hat, so glaube ich, dass ich wohl Freude an der Arbeit haben k"onnte, und dann w"urde ich mich auch recht gern einer Erweiterung der Triangulation nach Westen, falls sie mir aufgetragen w"urde, unterziehen. Die Anschliessung an die Krayenhoffschen Dreiecke ist allerdings w"unschenswerth, allein wo sind denn diese zu finden? Ich weiss nicht, ob sie irgendwo gedruckt sind, und der schlechte Erfolg mit den Eimbeck'schen Dreiecken macht mich ganz muthlos. Auch Laplace, an den ich vor etwa 9 Wochen geschrieben habe, hat mir gar nicht geantwortet.

Meiner Meinung nach sollten alle gut gemessenen Dreiecke 1.\ Ordnung als etwas betrachtet werden, worauf das ganze Publicum Anspruch hat, und nach und nach sollte ganz Europa mit solchen Dreiecken "uberzogen werden. Ich habe mir schon seit Jahren eine eigene Methode entworfen, wie solche Messungen am zweckm"assigsten behandelt werden k"onnen; denn alles, was ich dar"uber gelesen habe, finde ich herzlich werthlos. So haben sich z.\ B.\ viele Mathematiker grosse M"uhe mit der Aufgabe gegeben, aus Abst"anden vom Meridian und Perpendikel die L"ange und Breite zu berechnen, mit R"ucksicht auf die elliptische Gestalt der Erde, w"ahrend, so viel ich weiss, niemand vorher gefragt hat:
\begin{enumerate}[label=\arabic*.]
\item wie denn jene Abst"ande, so verstanden, wie man sie gew"ohnlich versteht, aus der Messung mit ebenso grosser Sch"arfe gefunden werden k"onnen; denn es scheint, dass die meisten diese Rechnung wie in der Ebene f"uhren, oder doch ganz unrichtige oder unbrauchbare Vorschriften daf"ur geben;
\item ob es denn "uberhaupt nur zweckm"assig sei, die so verstandenen Abst"ande zu gebrauchen, da es entschieden ist, dass, wenn man sie hinl"anglich scharf aus den Dreiecken ableiten will, dies nur durch h"ochst beschwerliche Rechnungen geschehen kann, so wie man aus ihnen nur mit vieler M"uhe wieder zu den L"angen und Breiten herabsteigt. Das Ganze w"urde nur ein >>die Pferde hinter den Wagen spannen<< sein.
\end{enumerate}
\emph{>>Soll etwas brauchbares zwischen die Dreiecke und die L"angen und Breiten gesetzt werden, so muss es etwas ganz anderes wie jene, so wie gew"ohnlich verstanden, Coordinaten sein.<<} Wie dies bei meiner Theorie geschieht, kann ich hier freilich nicht umst"andlich ausf"uhren; nur so viel bemerke ich, dass das, was ich zwischen die Dreiecke und die L"angen und Breiten setze, diejenigen Coordinaten sind, 1)~\emph{mit denen am zweckm"assigsten jeder Punkt in einer Ebene dargestellt werden kann.} Diese Coordinaten folgen h"ochst bequem und leicht aus den gemessenen Dreiecken, und ohne eine sehr genaue Kenntniss der Abplattung der Erde vorauszusetzen, und 2)~aus ihnen folgt wieder ebenso leicht die L"ange und Breite, nat"urlich indem man die Abplattung kennen muss. Ich habe die Absicht, diese Theorie, wo nicht fr"uher, doch mit meinen k"unftigen Messungen bekannt zu machen und bitte vorerst, diese angedeuteten Ideen noch f"ur sich zu behalten. Sehr gern w"urde ich sie nicht bloss auf die hannoverschen Dreiecke, sondern auf alle andern damit in Verbindung kommenden anwenden und so eine \emph{Description g\'{e}om\'{e}trique} eines grossen Theils von Europa geben, wenn ich durch Mittheilung geh"orig unterst"utzt w"urde. Aber!!

Herr v.~M"uffling hat mir doch seine 15 Dreiecke vom Rhein bis Seeberg mitgetheilt. Vorl"aufig, aber freilich nur sehr roh, habe ich bereits G"ottingen angeschlossen. Nemlich 1812 habe ich auf dem Hanstein, dessen Lage gegen G"ottingen n"aherungsweise aus meinen Winkelmessungen in hiesiger Gegend folgt, die Winkel zwischen G"ottingen, Brocken und der Boineburg (2.\ M"urring'scher Punkt) gemessen, freilich auf mehrere Minuten ungewiss; doch glaube ich, dass die Eintragung G"ottingens, die hieraus folgt, wohl auf 100 Meter beinahe zuverl"assig ist. Es folgt daraus: L"angenunterschied zwischen der G"ottinger und Seeberger Sternwarte in Zeit $3^{\mathrm{s}}\!\!\!,\!87$, was sehr nahe mit den astronomischen Bestimmungen zutrifft. Paris w"are hienach, wenn es $33^{\circ}35'$ westlich von Seeberg liegt, $30^{\mathrm{m}}\!\!\!,\!72633$ westlich von G"ottingen. Die neue Sternwarte liegt $139''$ "ostlich von der alten, die ich fr"uher immer $30^{\mathrm{m}}23\frac{1}{2}''$ von Paris setzte.

\medskip
\noindent \textbf{Gauss an Olbers.} G"ottingen, 18.~April 1822.

\medskip
\noindent Bei allen meinen Rechnungen liegen folgende Dimensionen der Erde zum Grunde
\begin{align*}
a &= 3272077.6 \quad \text{[Toisen]} \\
b &= 3261467.2 \quad \text{[Toisen]}
\end{align*}
\[\text{Abpl.} = \frac{1}{302.78}\,\footnote{Siehe die Briefe an Olbers vom 6.~Julius 1824, S.~320 Anmerkung, und vom 1.~M"arz 1827, S.~378.}\]
Eigentlich hatte ich ganz Walbeck's Resultat annehmen wollen; durch einen Schreibfehler hatte ich aber jene schon vor l"angerer Zeit der Berechnung von mancherlei H"ulfstafeln untergelegt, und hielt es um so weniger der M"uhe werth, diese deshalb umzuarbeiten, da der Unterschied weit unter der durch alle Gradmessungen zu erreichenden Genauigkeit liegt.

\newpage
\noindent \textbf{Gauss an Olbers.} Zeven, 4.~Julius 1824.

\medskip
\noindent F"ur weiteres Fortschreiten sind die Aussichten ausserst schlecht. Nach meiner fr"uhern Hoffnung Bremen zu umgehen, wird nicht thunlich sein, da der Weierberg mit Brittendorf nicht zu verbinden ist. Aber auch gar nichts anderes rechtliches l"asst sich mit Brittendorf im Nordwesten oder Norden verbinden. M"obius' Recognoscirung von Bremerv"orde bis Osterholz hat durchaus gar kein Resultat gegeben. Der einzige Punkt w"are bei Wentel, der aber an Bremen wohl nur einen Winkel von etwa $12^{\circ}$, an Brittendorf einen von $46^{\circ}$, an Wentel von $122^{\circ}$ \emph{praeterpropter} geben w"urde, und wo ich auch noch gar nicht weiss, ob das Opfer, was ich durch ein so schlechtes Dreieck br"achte (und welches etwas gebessert w"urde, wenn sich Wentel zugleich mit Bottel verbinden liesse), durch eine einigermassen rechtliche Aussicht nach Nordwesten von Wentel aus compensirt w"urde. Am Ende werde ich also doch vielleicht den Steinberg noch mit zuziehen m"ussen, um mich s"ud"ostlich um Bremen herum zu drehen, oder Bremen nur wie eine vorgeschobene Zunge betrachten und die Verbindung mit Krayenhoff n"ordlich "uber Stade, oder s"udlich "uber Osnabr"uck suchen m"ussen. Das erstere allein zu thun, ist wegen der unerh"ort schlechten Beschaffenheit der nord"ostlichen Krayenhoffschen Dreiecke (wor"uber ich Ihnen fr"uher einmal geschrieben) auch wohl bedenklich. Sehr w"unschte ich Ihre Ansicht dar"uber zu haben\,\ldots\ldots\ldots

\medskip
\noindent \textbf{Gauss an Olbers.} Zeven, 8.~Julius 1824.

\medskip
\noindent Ihre beiden letzten g"utigen Briefe habe ich richtig erhalten\,\ldots\ldots\ldots

Der erstere hat mich r"ucksichtlich aller meiner Messungen sehr niedergeschlagen. Da Sie das Dreieck Bremen-Br"uttendorf-Wentel wegen des zu spitzen Winkels an Bremen verwerfen, so brechen Sie dadurch zugleich den Stab "uber die, wie es scheint, einzig m"ogliche Art, auf der andern Seite um Bremen herum zu kommen; denn in dem Dreieck Bremen-Bottel-Steinberg wird der Winkel in Bremen noch viel spitzer sein. Sie setzen zwar mit Ihrer gewohnten G"ute hinzu, dass doch bei jenem Dreieck die Genauigkeit desswegen nicht bedeutend leiden w"urde, weil ich in meine Messungen eine so grosse Sch"arfe lege. Allein dieser Grund, dessen Wahrheit ich jetzt auf sich beruhen lassen will, kann mich durchaus im geringsten nicht beruhigen.

\newpage
\noindent Nach meinem Grundsatze soll man immer, so genau man nur kann, beobachten; der Grad der Genauigkeit in den Beobachtungen, gleichviel wie gross oder wie klein er sein mag, bedingt immer wieder den Grad der Genauigkeit, die man von den Resultaten fordern darf, und die Genauigkeit der Beobachtungen kann nach meiner Meinung ein an sich schlechtes Dreieck durchaus nicht gut machen; wenigstens w"are sonst "uberfl"ussig gewesen, die "ubrigen guten Dreiecke mit derselben Sch"arfe zu messen. H"ochstens kann dadurch dann das ganze System wieder in Parallele mit andern an sich viel schlechtern Messungen zur"uckkommen.

Ich habe es bisher f"ur ein blosses Vorurtheil gehalten, wenn man Dreiecke mit sehr kleinen Winkeln der Genauigkeit f"ur nachtheilig hielt, insofern die den spitzen Winkeln gegen"uberliegenden Seiten keine "Ubergangsseiten abgeben; ich habe solche kleinwinkelige Dreiecke bloss desswegen f"ur minder gut gehalten, weil man damit auf einmal nicht viel weiter kommt, also mehr Zeit und Kosten gebraucht, als wenn man auf einmal viel fortschreiten kann; und auch dieser Grund f"allt ganz weg, wenn die Aufsuchung und Instandsetzung eines grossen Dreiecks vielleicht doppelt so viel Zeit kostet, als die Messung zweier Dreiecke zusammen, die eben dahin f"uhren, und wovon das eine einen sehr spitzen Winkel hat. Demungeachtet habe ich nicht ganz nach diesem Princip gehandelt, sondern ein Dreieck mit einem kleinen Winkel nie eher adoptirt, als bis ich fast alle M"oglichkeiten ersch"opft hatte, es zu vermeiden (nur diejenige M"oglichkeit nicht, die zu schlechten Messungen selbst gef"uhrt hatte, d.\ i.\ Zach's hohe Th"urme); nicht weil ich geglaubt h"atte, dadurch an Genauigkeit etwas zu gewinnen, sondern aus dem wohl verzeihlichen Wunsche, dem System so viel m"oglich, ausser dem innern Gehalt, auch Sch"onheit und R"undung zu geben.

Da ich nun aber Sie durch das, was ich in einem fr"uhern Briefe dar"uber schrieb, nicht "uberzeugt habe, sondern da Sie den Nachtheil, der f"ur die Genauigkeit aus dem spitzen Winkel sonst entstehen w"urde, durch die Sch"arfe der Messungen gut gemacht verlangen, was nach meiner Ansicht unm"oglich ist, so werde ich selbst in meiner bisherigen Ansicht ganz irre und zweifelhaft, ob sie nicht ganz unrichtig gewesen, und darf wenigstens auf keinen Fall hoffen, andere von der Richtigkeit derselben zu "uberzeugen. Was namentlich das Dreieck Bremen-Br"uttendorf-Wentel betrifft, so hatte ich mich selbst sehr ungern dazu entschlossen, weil es nicht sch"on ist, und auf einmal nicht viel weiter bringt; r"ucksichtlich der Genauigkeit aber (ganz abgesehen davon, wie genau die Winkelmessungen an sich sind), w"urde ich dasselbe, seine Winkel zu $12^{\circ}$, $46^{\circ}$, $122^{\circ}$ angenommen, vollkommen einem andern gleichgestellt haben, dessen Winkel $76^{\circ}$, $46^{\circ}$, $58^{\circ}$ gewesen w"aren.

\newpage
\noindent \textbf{Gauss an Olbers.} G"ottingen, 19.~Februar 1825.

\medskip
\noindent Ich schicke Ihnen mein vollst"andiges H"ohenverzeichniss\footnote{In Gauss' Nachlass ist das Original dieses Verzeichnisses vorhanden; demselben ist das folgende zugef"ugt:
\begin{quote}
>>Die H"ohen bei den Hauptpunkten beziehen sich, mit Ausnahme von G"ottingen, L"uneburg und Hamburg, wo das N"ahere besonders bemerkt ist, allemal auf diejenige Fl"ache, auf welcher der Repetitionskreis gestanden hat. Bei den meisten Pl"atzen ist dies ein aufgemauertes steinernes Postament von $3\frac{1}{2}$ Fuss H"ohe "uber der Erde, zuweilen ein paar Zoll mehr, zuweilen weniger. Nur beim Meridianzeichen betr"agt diese H"ohe etwas mehr, nemlich 5 Fuss.<<
\end{quote}}. Bei Hamburg ist eine kleine Ver"anderung von $+2,7$ Fuss vorgenommen, da Schumacher mir den H"ohenunterschied zwischen dem Knopf und den Fenstern des Cabinets, den ich fr"uher zu 47,9 Fuss nach Benzenberg's Kupfer(stich) angenommen hatte, nach Sonnin's Kupfer[stich] zu 53,3 angibt; ich habe also einstweilen den Knopf (meinen Zielpunkt) um 2,7 h"oher, die Fenster 2,7 [Fuss] tiefer gesetzt als vorher. Dadurch wird dann auch Schumacher's Barometer 2,7 [Fuss] tiefer als vorher, also gegen G"ottingen $-357,8$ [Fuss] und gegen Ihr Barometer $+70,4$ [Fuss] $= 11,73$ Toisen. Der Unterschied mit Ihrer Barometerbestimmung ist also jetzt ganz unbedeutend. Doch wird dies noch etwas modificirt werden, da ich auch auf einigen Stationen den Fussboden der Laterne zum Zielpunkt gebraucht habe, dessen Tiefe unter dem Knopfe mir Schumacher nicht mitgetheilt hat, und "uberhaupt sollte wohl die relative H"ohe der drei Punkte ordentlich trigonometrisch gemessen werden (aus einem nahen Standpunkte).

\newpage
\noindent \textbf{Relative H"ohen.}\\[1ex]
\begin{tabular}{r|l|c}
\multicolumn{3}{c}{\textit{Pariser Fuss}} \\[0.5ex]
\hline
1 & G"ottingen, Sternwarte, Fussboden & 0 \\
2 & N"ordl.\ Meridianzeichen (Oberfl.\ des Postaments) & $+201,\!0$ \\
3 & Hohehagen (Postamentsoberfl"ache) & $+1072,\!6$ \\
4 & Hils (P.) & $+841,\!2$ \\
5 & Brockenhaus (Marmortisch auf dem Thurm) & $+3061,\!7$ \\
6 & Lichtenberg (P.) & $+274,\!1$ \\
7 & Deister, Calenberg (P.) & $+468,\!2$ \\
8 & Garssen (P.) & $-242,\!5$ \\
9 & Falkenberg (P.) & $-15,\!5$ \\
10 & Scharnhorst (P.) & $-190,\!6$ \\
11 & Breithorn (P.) & $-109,\!0$ \\
12 & Hauselberg (P.) & $-109,\!2$ \\
13 & Wulfsode (P.) & $-158,\!5$ \\
14 & Timpenberg (P.) & $-120,\!1$ \\
15 & Nindorf (P.) & $-120,\!0$ \\
16 & L"uneburg, Michaelis, Fussboden der Laterne & $-243,\!2$ \\
17 & \quad\quad\quad Knopf & $-178,\!6$ \\
18 & \quad\quad\quad Spitze & $-168,\!3$ \\
19 & \quad\quad\quad Johannis, Knopf & $-109,\!5$ \\
20 & \quad\quad\quad \quad\quad\quad Spitze & $-101,\!4$ \\
21 & \quad\quad\quad Nicolai, Knopf & $-199,\!3$ \\
22 & \quad\quad\quad \quad\quad\quad Spitze & $-188,\!4$ \\
23 & \quad\quad\quad Lamberti, Knopf & $-215,\!7$ \\
24 & \quad\quad\quad \quad\quad\quad Spitze & $-207,\!7$ \\
25 & \quad\quad\quad H"ochste Stelle des Kalkberges & $-291,\!6$ \\
26 & \quad\quad\quad Platz nahe vor dem N.W.\ Thor & $-393,\!1$ \\
27 & L"une, Klosterthurm, Knopf & $-331,\!2$ \\
28 & \quad\quad\quad\quad\quad\quad Spitze & $-325,\!0$ \\
\hline
\end{tabular}

\newpage
\begin{tabular}{r|l|c}
\hline
29 & Hamburg, Michaelisthurm, Knopf & $-31,\!2$ \\
30 & \quad\quad\quad\quad\quad\quad\quad\quad\quad\quad Fenster des ob.\ Cab. & $-19,\!6$ \\
31 & Wilsede (P.) & $+86,\!5$ \\
32 & Elmhorst (P.) & $+96,\!2$ \\
33 & Litberg (P.) & $+99,\!4$ \\
34 & Bullerberg (P.) & $+86,\!5$ \\
35 & Bottel (P.) & $+21,\!0$ \\
36 & Steinberg (P.) & $+50,\!7$ \\
37 & Br"uttendorf (P.) & $+48,\!7$ \\
38 & Zeven, Kirchthurm, Fussboden der Laterne & $-14,\!1$ \\
39 & \quad\quad\quad\quad\quad\quad\quad\quad Knopf & $-10,\!0$ \\
40 & \quad\quad\quad\quad\quad\quad\quad\quad Spitze & $-1,\!6$ \\
41 & \quad\quad\quad\quad\quad\quad\quad Garten beim Posthause, nahe der Aue & $-9,\!1$ \\
42 & \quad\quad\quad\quad\quad\quad\quad Platz auf dem Felde beim s"udlichen Eingang des Dorfs & $-5,\!8$ \\
43 & Bremen, Ansgarius, Fussboden der Laterne & $-2,\!4$ \\
44 & \quad\quad\quad\quad\quad\quad\quad\quad\quad\quad\quad Knopf & $+5,\!5$ \\
45 & \quad\quad\quad\quad\quad\quad\quad\quad\quad\quad\quad Spitze & $+11,\!5$ \\
46 & \quad\quad\quad\quad\quad\quad\quad\quad\quad\quad\quad Liebenfrauen, Knopf & $+9,\!0$ \\
47 & \quad\quad\quad\quad\quad\quad\quad\quad\quad\quad\quad Dom, Knopf & $+17,\!8$ \\
48 & \quad\quad\quad\quad\quad\quad\quad\quad\quad\quad\quad \quad\quad\quad\quad Spitze & $+26,\!1$ \\
49 & \quad\quad\quad\quad\quad\quad\quad\quad\quad\quad\quad Martini, Knopf & $+7,\!9$ \\
50 & \quad\quad\quad\quad\quad\quad\quad\quad\quad\quad\quad Gymnasium, Knopf & $+21,\!2$ \\
51 & \quad\quad\quad\quad\quad\quad\quad\quad\quad\quad\quad Strassenpflaster am Domhof unweit des alten Museums & $-0,\!3$ \\
52 & \quad\quad\quad\quad\quad\quad\quad\quad\quad\quad\quad Windm"uhlenberg am Herdner Thor & $+20,\!4$ \\
53 & \quad\quad\quad\quad\quad\quad\quad\quad\quad\quad\quad Garten beim Hause des Hrn.\ Dr.\ Olbers & $-4,\!0$ \\
54 & \quad\quad\quad\quad\quad\quad\quad\quad\quad\quad\quad Fussboden in Hrn.\ Dr.\ Olbers Observ. & $-4,\!2$ \\
55 & \quad\quad\quad\quad\quad\quad\quad\quad\quad\quad\quad Barometer-Gef"ass, daselbst & $-4,\!6$ \\
56 & Brillit (P.) & $-114,\!1$ \\
57 & Garlster Haide (P.) & $+44,\!5$ \\
\hline
\end{tabular}

\newpage
\noindent \textbf{Gauss an Olbers.} G"ottingen, 25.~Februar 1825.

\medskip
\noindent Jetzt noch ein paar Worte "uber das geod"atische Nivellement.

Ich bin zwar selbst mit mir "uber diesen Gegenstand ganz auf dem Reinen, allein ich weiss nicht, ob ich mich in der K"urze so dar"uber werde erkl"aren k"onnen, dass ich Sie sofort zur Uebereinstimmung bringen werde.

Ich habe immer geglaubt, dass der Ausdruck \emph{>>Localattraction<<} sehr "ubel gew"ahlt ist und leicht verkehrte Ansichten veranlassen kann. Man sollte sagen, dass die Richtungen der Schwere nicht mit dem Gange, der bei einem gleichf"ormigen Sph"aroid stattfinden w"urde, Schritt halten. Die Richtung der Schwere ist das Totalproduct der Anziehung aller Bestandtheile des Erdk"orpers (und der Centrifugalkraft), und bei dessen unregelm"assiger Zusammensetzung in R"ucksicht der Dichtigkeit, sowie bei den Unebenheiten auf der Oberfl"ache wird jene nicht dieselbe sein k"onnen, wie bei einem regelm"assigen Sph"aroid. Allein wie auch die Zusammensetzung sei, immer wird durch jeden Punkt eine Fl"ache, die ganz um die Erde herum geht, gelegt werden k"onnen, auf welcher die Richtung der Schwere genau senkrecht ist, und die Oberfl"ache einer zusammenh"angenden ruhigen Fl"ussigkeit w"urde dieselbe vorstellen. Diese Fl"ache ist es, die eine Horizontalf"ache heisst (\emph{couche de niveau}); den Punkten dieser Fl"ache legt man gleiche H"ohe bei, ohne sich im mindesten darum zu bek"ummern, ob oder wie viel sie von einem elliptischen Sph"aroid abweichen, und die H"ohen "uber dieser Fl"ache gibt sowohl das Barometer als die trigonometrische Messung an, so dass beide immer mit einander "ubereinstimmen m"ussen. Dabei wird bloss vorausgesetzt\footnote{und nat"urlich auch, dass alle Zenithdistanzen reciprok gemessen werden, was bei meiner Messung ohne Ausnahme gilt. Bei einseitigen Messungen ist Ihre Bemerkung vollkommen gegr"undet, da man dabei die Amplitude sph"aroidisch berechnen muss.}, dass auf jeder Dreieckslinie die Richtung der Schwere sich nach dem Gesetz der Stetigkeit "andert (obgleich vielleicht schneller oder langsamer als bei dem elliptischen Sph"aroid), und diese Voraussetzung kann nur dann eine kleine Unrichtigkeit hervorbringen, wenn an der einen Dreiecksstation eine wahre Localattraction stattfindet, die bloss "ortlich und auf einen kleinen Raum beschr"ankt (ausserhalb desselben unmerklich) ist. Allein ich halte mich "uberzeugt, dass, den Brocken h"ochstens ausgenommen, eine solche Localattraction im ganzen Umfange meiner und der Schumacherschen Dreiecke nicht statthat.\,\ldots\ldots\ldots

\newpage
\noindent \textbf{Gauss an Olbers.} G"ottingen, 9.~October 1825.

\medskip
\noindent Ich habe dieser Tage angefangen, in Beziehung auf mein k"unftiges Werk "uber H"ohere Geod"asie, einen (sehr) kleinen Theil dessen, was die krummen Fl"achen betrifft, in Gedanken etwas zu ordnen. Allein ich "uberzeuge mich, dass ich bei der Eigenth"umlichkeit meiner ganzen Behandlung des Zusammenhanges wegen gezwungen bin, sehr weit auszuholen, so dass ich sogar meine Ansicht "uber die Kr"ummungshalbmesser bei planen Curven vorausschicken muss. Ich bin dar"uber fast zweifelhaft geworden, ob es nicht gerathener sein wird, einen Theil dieser Lehren, der ganz rein geometrisch (in analytischer Form) ist und Neues mit Bekanntem gemischt in neuer Form enth"alt, erst besonders auszuarbeiten, ihn vielleicht von dem Werke abzutrennen und als eine oder zwei Abhandlungen in unsere Commentationen einzur"ucken. Indessen kann ich noch vorerst die Form der Bekanntmachung auf sich beruhen lassen, und werde einstweilen in dem zu Papier bringen fortfahren.

\medskip
\noindent \textbf{Gauss an Olbers.} G"ottingen, 2.~April 1826.

\medskip
\noindent R"ucksichtlich meiner Messungen kann ich f"ur mich allein wenig beschliessen. Mein Auftrag ist im Grunde r"ucksichtlich des trigonometrischen Theils vollendet, und ich dachte, insofern Schumacher mich geh"orig unterst"utzen will, im Sp"atsommer die Zenithsector-Beobachtungen vorzunehmen. Der wankende Zustand meiner Gesundheit schreckt mich ab, auf erweiterte trigonometrische Messungen anzutragen; inzwischen habe ich vor kurzem, unter uns gesagt, in einem Schreiben an M"unster erkl"art, dass ich bereit bin, meine Kr"afte auch noch k"unftig darauf zu wenden, falls solche gefordert werden sollten.

Meine theoretischen Arbeiten lassen bei ihrem so sehr grossen Umfange leider noch viele L"ucken; am leichtesten w"are mir geholfen, wenn ich mir erlaubte, mit der Bekanntmachung meiner Messungen zwar alle meine Rechnungseinrichtungen zu verbinden, aber deren Ableitung aus ihren h"ohern Gr"unden f"ur ein ganz getrenntes Werk f"ur gl"ucklichere zuk"unftige Zeiten aufzusparen. Dann w"are nirgends ein Anstoss. Vors erste werde ich die scharfe Ausgleichung meiner 32 Punkte, die 51 Dreiecke und 146 Richtungen liefern, vornehmen. Die H"ohenausgleichung (ein sehr viel leichteres Gesch"aft) habe ich in diesen Tagen vollendet.\,\ldots\ldots\ldots

\newpage
\noindent \textbf{Gauss an Olbers.} G"ottingen, 14.~Januar 1827.

\medskip
\noindent Ich glaube Ihnen schon fr"uher gemeldet zu haben, dass ich es ganz unthunlich gefunden habe, dasjenige Werk, welches ich "uber meine Messungen in Zukunft zu geben denke, auch in theoretischer R"ucksicht ganz selbstst"andig zu machen, wenn ich nicht wenigstens einen grossen Theil des Theoretischen vorher anderswo besonders behandle. Es wird schon volumin"os, die Gr"unde meiner Operationen zu entwickeln, aber, wenn ich mich so ausdr"ucken darf, die Gr"unde der Gr"unde k"onnen nicht in das Werk selbst kommen, ohne es ganz buntscheckig und doch unbefriedigend zu machen. Ich habe mich daher entschlossen, verschiedene theoretische Materien erst abgesondert in einzelnen Abhandlungen zu entwickeln, wodurch es auch allein m"oglich wird, diese bedeutenden neuen Capitel der Mathematik mit einem gewissen Grade von Vollst"andigkeit auszuf"uhren. Gewissermassen ist meine Preisschrift "uber die Transformation der Fl"achen die erste dieser Abhandlungen; die zweite habe ich vor einigen Monaten der k.\ Societ"at "ubergegeben als \emph{Supplementum theoriae combinationis observationum etc.} Sie enth"alt die Principien, die als Grundlage der Ausgleichung der Beobachtungen angewandt werden m"ussen, und selbst einige Beispiele von Krayenhoff's und meinen Messungen.

\medskip
\noindent \textbf{Gauss an Olbers.} G"ottingen, 1.~M"arz 1827.

\medskip
\noindent Meine Abhandlung, oder vielleicht richtiger, meine erste Abhandlung "uber die krummen Fl"achen habe ich vollendet; ich werde sie aber der Societ"at noch nicht "ubergeben, da doch auf die Ostermesse kein Band herauskommt. Die beiden von mir 1825 und 1826 "ubergebenen Abhandlungen "uber die biquadratischen Reste und \emph{Suppl.\ theor.\ combin.\ observ.} sind noch nicht zu drucken angefangen. Jene Abhandlung enth"alt zur unmittelbaren Benutzung in meinem k"unftigen Werke "uber die Messung eigentlich nur ein paar S"atze, nemlich 1)~was zur Berechnung des Excesses der Summe der 3 Winkel "uber $180^{\circ}$ in einem Dreiecke auf einer nicht sph"arischen Fl"ache, wo die Seiten k"urzeste Linien sind, erforderlich ist, 2)~wie in diesem Fall der Excess auf die drei Winkel ungleich vertheilt werden muss, damit die Sinus den Seiten gegen"uber proportional werden. In praktischer R"ucksicht ist dies zwar ganz unwichtig, weil in der That bei den gr"ossten Dreiecken, die sich auf der Erde messen lassen, diese Ungleichheit in der Vertheilung unmerklich wird; aber die W"urde der Wissenschaft erfordert doch, dass man die Natur dieser Ungleichheit klar begreife. Und so kann man allerdings hier, wie "ofters, ausrufen: \emph{Tantae molis erat!} um dahin zu gelangen. --- Wichtiger aber als die Aufl"osung dieser 2 Aufgaben ist es, dass die Abhandlung mehrere allgemeine Principien begr"undet, aus denen k"unftig, in einer speciellern Untersuchung, die Aufl"osung von einer Menge wichtiger Aufgaben abgeleitet werden kann.

\newpage
\noindent Ich komme noch einmal auf den im Anfange dieses Briefes erw"ahnten Gegenstand zur"uck. Ich habe bei allen meinen H"ulfstafeln und Rechnungen Walbeck's Abplattung $\frac{1}{302.78}$ zum Grunde gelegt; aber ich glaube, dass die s"ammtlichen bisherigen Gradmessungen, wenn man ihre Data vollst"andig aufn"ahme, zeigen w"urden, dass die Schumacher'sche Abplattung $\frac{1}{331.7}$ sich beinahe ebenso gut damit w"urde vereinigen lassen, und es scheint mir, dass alle bisher vorhandenen Gradmessungen noch viel zu kleine Ausdehnung haben, um die Abplattung in engere Grenzen einzuschliessen. Was gewiss ist, ist, dass die Erde ein unregelm"assiger K"orper ist; die Polh"ohen der Orter (ganz abstrahirt von Beobachtungsfehlern) schwanken immer mehrere Secunden (vielleicht hie und da ziemlich viele Secunden) um die Werthe, die man unter Voraussetzung irgend eines regelm"assigen Ellipsoids berechnet, und ebenso schwanken die wirklichen Pendell"angen um die berechneten (das mittlere Schwanken der Pendell"angen vielleicht 4 englische Linie). Aber eben deshalb k"onnen Gradmessungen von kleiner Ausdehnung wenig zur Kenntniss der mittlern Gestalt der Erde beitragen, namentlich halte ich den lappl"andischen Bogen f"ur viel zu klein. Um so wichtiger, daucht mir, ist es, dass nach und nach alle scharfen Triangulirungen in Europa in Einen Zusammenhang gebracht werden. Ich habe jetzt Hoffnung, die bayerischen Dreiecke mitgetheilt zu erhalten. Wenn erst alle europ"aischen Sternwarten von Abo bis Palermo und von Nicolajef bis Dublin durch Dreiecke zusammenh"angen, so dass ihre relativen Lagen gegen einander mit aller Genauigkeit, die die feinsten geod"atischen Messungen verschaffen k"onnen, bestimmt sind, so wird man, d.\ i.\ so werden unsere Nachkommen, alles mit viel mehr Sicherheit beurtheilen k"onnen.\,\ldots\ldots\ldots

\newpage
\noindent \textbf{Gauss an Olbers.} G"ottingen, 14.~Junius 1830.

\medskip
\noindent Die Zeit, die mir in diesem Sommer zu eigener Arbeit "ubrig bleibt, denke ich der Fortsetzung meiner Abhandlung "uber die biquadratischen Reste zu widmen, damit diese Arbeit, deren erste Anf"ange sich schon von 1805 her datiren, ihrer Vollendung n"aher komme. Sie wird wenigstens noch zwei ausgedehnte Abhandlungen erfordern; vorerst werde ich es aber bei einer bewenden lassen, und die erste mir dann wieder zu Theil werdende M"usse erst wieder einem andern Gegenstande widmen, wahrscheinlich den theoretischen Methoden der h"ohern Geod"asie. Es ist seit geraumer Zeit mein Schicksal gewesen, immer solche Arbeiten aufzunehmen, bei denen sich in der Darstellung nicht schnell fortschreiten l"asst.\,\ldots\ldots\ldots

\medskip
\noindent \textbf{Gauss an Olbers.} G"ottingen, 2.~September 1837.

\medskip
\noindent Gerling, der seine trigonometrischen Messungen in Hessen jetzt beendigt hat, hat jetzt noch eine Operation veranstaltet, die zur Bestimmung des L"angenunterschiedes zwischen G"ottingen und Mannheim dienen soll. Es werden Signale auf zwei Bergen, Meisner und Feldberg, gegeben; erstere sind in meiner, letztere in der Mannheimer Sternwarte sichtbar; beide aber zugleich auf einem Zwischenberge bei Marburg, wo Gerling mit einem Kessler'schen Chronometer beobachtet. Die Signale sind Pulverzeichen bei Nacht und heliotropische bei Tage. Das Wetter ist nicht g"unstig; hier sind zwar bisher schon ziemlich viele Zeichen von beiderlei Art beobachtet, aber Gerling hatte nach seinem letzten Briefe noch fast nichts vom Feldberge her gesehen.\,\ldots\ldots\ldots

\newpage
\subsection*{[5.]}
\addcontentsline{toc}{subsection}{Gauss an Gerling}

\noindent \textbf{Gauss an Gerling.} G"ottingen, 5.~October 1821.

\medskip
\noindent Meinen herzlichen Gl"uckwunsch zu dem von Ihnen "ubernommenen Gesch"aft. Es freut mich, dass es in Ihre H"ande kommt, da Sie es sich zur Pflicht machen werden, das ganze Triangelnetz mit aller m"oglichen Sorgfalt auszuf"uhren. Ich halte dies f"ur etwas "uberaus w"unschenswerthes, und beklage es immer, wenn man schon bei den Triangeln der ersten Ordnung geizig "uberlegt, durchaus nichts mehr an Genauigkeit anwenden zu wollen, als f"ur den allerletzten Zweck unumg"anglich n"othig ist. Die genaueste Kenntniss der relativen Lagen der interessantesten Punkte eines Landes kann in vielfacher Beziehung n"utzlich sein, auch ganz abgesehen davon, dass eine Detailvermessung darauf am besten zu st"utzen ist. Es w"are gewiss ausserst wichtig, wenn der gr"osste Theil von Europa vollst"andig mit Einem Netz "uberzogen w"are, und nach und nach werden wir dahin kommen; jeder Staat sollte es sich zur Ehre rechnen, seinen Antheil daran so gut zu liefern, dass er w"urdig sei, neben den besten zu stehen.

Ich sollte glauben, dass Sie gut thun w"urden, den Meisner auch mit zu einem Dreieckspunkt zu w"ahlen; Sie werden denselben unmittelbar mit wenigstens 5 meiner Dreieckspunkte in Verbindung setzen k"onnen, und wahrscheinlich mit ebenso viel M"uffling'schen dazu. Ich werde gern dabei hilfreiche Hand bieten, und insofern es sich nur irgend einrichten l"asst, die diesseitigen Beobachtungen dahin machen. Die Heliotrope erleichtern solche Arbeiten ausserordentlich, und es bedarf auf dem Meisner weiter keiner Anlage, als dass ein steinernes Postament von circa $1\frac{1}{2}$ Fuss Quadrat und 3 Fuss H"ohe "uber der Erde gesetzt werde, um Heliotrop und Theodolithen darauf zu stellen, und dass vielleicht einige B"aume gef"allt werden, falls ohne das die freie Aussicht nicht vollst"andig erreicht werden kann.\,\ldots\ldots\ldots

\newpage
\noindent \textbf{Gauss an Gerling.} G"ottingen, 21.~Februar 1822.

\medskip
\noindent Wenn der Anschluss an die G"ottinger Sternwarte und der an den Inselsberg nicht auf Einem Platze auf dem Meisner geschehen kann, so d"urfen Sie nun beide Pl"atze anwenden. Es ist dabei gar nicht viel Vermehrung der Arbeit; denn nach den Principien der Wahrscheinlichkeitsrechnung brauchen Sie die Winkel zwischen den Objecten, die an beiden Punkten sichtbar sind, insofern Sie sie an beiden messen, an jedem nur halb so oft zu repetiren, als Sie gethan haben w"urden, wenn Sie nur Einen Standpunkt gebraucht h"atten. Die gegenseitige Lage beider Punkte gegen einander mit der gr"ossten Sch"arfe und doch ohne "uberfl"ussige M"uhe zu erhalten, wird Ihnen nicht schwer fallen.\,\ldots\ldots\ldots

\medskip
\noindent \textbf{Gauss an Gerling.} G"ottingen, 7.~November 1822.

\medskip
\noindent Bei sehr flacher Incidenz, wo auch zuletzt die M"oglichkeit der Lenkung aufh"ort, habe ich immer mit dem herrlichsten Erfolg doppelte Reflexion anwenden lassen, indem nicht die Sonne selbst, sondern ein auf der Erde an schicklicher Stelle aufgestellter Handspiegel den Heliotrop speiste. Ich empfehle Ihnen diesen Kunstgriff zur Nachahmung.\,\ldots\ldots\ldots

In einem bergigen Lande ist es eine Lust zu messen, desto gr"osser, je gr"osser die Entfernungen sind, die beim Gebrauch der Heliotropen gar keine Grenzen haben. Mein grosser Spiegel, 1 Quadratfuss, war sogar durch den dichten Moorbrandsqualm, der mehrere Quadratmeilen bedeckte, von Lichtenberg nach Falkenberg, fast 12 Meilen, durchgedrungen, welcher Qualm freilich f"ur die winzigen Heliotropspiegel zu dicht gewesen war.\,\ldots\ldots\ldots

\newpage
\noindent \textbf{Gauss an Gerling.} G"ottingen, 27.~Julius 1823.

\medskip
\noindent Ich habe mir von Repsold noch ein neues Ocular f"ur meinen Theodolithen verfertigen lassen, mit Spinnenf"aden von fast unglaublicher Feinheit, $29''$ von einander abstehend. (Die in dem Ertel'schen Ocular, auch schon recht fein, aber viel dicker als jene, sind $36''$ von einander.) Bei Beobachtung entfernter blasser irdischer Gegenst"ande ist jene Einrichtung "ausserst vortheilhaft; leider bekam ich sie nur erst \emph{post festum}. Der in Nindorf, Timpenberg und L"uneburg fast immer bei dem herrauchigen Zustand der Luft sehr blass aussehende Hamburger Thurm hat mich sehr geplagt, ebenso wie der L"uneburger in Hamburg, und im allgemeinen sind die Messungen in jenen Gegenden nicht ganz so scharf wie die fr"uhern, wozu "ubrigens auch das ewige Schwanken der Th"urme mit beigetragen hat. Der Michaelis-Thurm in Hamburg ist, so lange ich dagewesen bin, nie ruhig gewesen; die horizontalen pendelartigen Schwankungen gehen oft "uber $\pm 1$ Minute. Dieser Thurm ist von allen meinen Dreieckspunkten als Standpunkt und (den Brocken abgerechnet) auch als Zielpunkt der allerschlechteste gewesen. Repsold verfertigt auch neue Libellen von ausserordentlicher Vollkommenheit, die er anstatt mit Weingeist mit Naphtha f"ullt. Sie kommen sehr viel schneller als die andern zur Ruhe.

\medskip
\noindent \textbf{Gauss an Gerling.} G"ottingen, 11.~August 1823.

\medskip
\noindent Ihr Ausdruck, wenn ich vom Brocken aus auf Heliotroplicht pointiren wolle und nicht das H"auschen vorz"uge etc., scheint auf einiges Missverst"andniss der Motive, um derenwillen ich zum zweiten Male auf den Brocken gehen k"onnte, hinzudeuten. Die scharfe Bestimmung der Richtung zum Inselsberg ist in der That ein Hauptzweck, und von Pointiren auf das Haus kann bei dieser Entfernung und der geringen H"ohe und Irregularit"at desselben gar keine Rede sein. Heliotroplicht vom Inselsberg zum Brocken ist also nat"urlich die \emph{Conditio sine qua non} meines Hingehens.

Dieser Eine Grund allein w"urde mich jedoch noch nicht wegen des abermaligen Besuchs rechtfertigen. Meine Winkelmessungen 1821 nach meinen eigenen Hauptpunkten waren bei dem ambulirenden Heliotrop und ung"unstigen Wetter zu d"urftig und des Ganzen nicht v"ollig w"urdig ausgefallen; daher ich diese durch gleichzeitige Besetzung mit Heliotropen wiederholen, wie nachher auch noch einmal zum Hohehagen und Hils zur"uckzukehren w"unschte.

\newpage
\noindent So wie ich nun mit gr"osstem Vergn"ugen vom Brocken und Hohehagen zum Inselsberg Licht f"ur Licht zu schicken erb"otig bin, so bin ich doch nicht im Stande, Ihnen einen meiner Heliotropen hiebei abzulassen; denn wenn ich Hohehagen, Hils und vielleicht Lichtenberg durch einen ambulirenden besetzen m"usste, so k"ame ich nicht weiter als 1821.

Es scheinen mir nun mehrere Wege, wie hier durchzukommen ist, denkbar:
\begin{enumerate}[label=\arabic*)]
\item Kommen wir gleichzeitig resp.\ zum Inselsberg und Brocken und schicken uns gegenseitig Licht, so wie mein Gehilfe Ordre haben soll, Ihnen von Zeit zu Zeit Licht vom Hohehagen, und der Ihrige mir vom Meisner zu schicken, so kann ich unter Beg"unstigung des Wetters, wie es Anfang September h"aufig ist, in wenigen Tagen, vielleicht in dreien, auf dem Brocken fertig werden, den Gehilfen mit dem Heliotrop dort lassen (der fortw"ahrend Ihnen zum Inselsberg Licht schickt, so wie der auf dem Hohehagen), zum Hohehagen eilen, Ihnen Zeichen meiner Ankunft geben und wieder in ein paar Tagen (unter Gunst des Wetters) dort das zum Inselsberg geh"orige absolviren, worauf Sie dann "uber Ihren Inselsberg-Heliotrop nach Gefallen disponiren k"onnen. Vom Hohehagen erhalten Sie dann best"andig Licht, so lange ich da bin, und wenn Sie es noch bed"urfen, ab und an, so bald ich auf dem Hils bin.

In Beziehung auf diesen Plan wird also alles darauf ankommen, ob Sie die Verbindung mit dem Brocken und den hannoverschen Dreiecken f"ur wichtig genug halten, die Besetzung des Kn"ulls und der Milseburg mit Heliotropen so lange aufzuschieben, bis jener effectuirt ist.
\item Obgleich der oben angezeigte Weg mir der liebste w"are, insofern dadurch die Verbindung der Messungen am vollkommensten effectuirt wird, so k"onnten Sie doch den Inselsberg-Heliotrop allenfalls schon fr"uher abgeben, nemlich so bald ich den Brocken verlasse, was durch Signale kundgemacht werden kann. In Beziehung auf die Richtung Hohehagen-Inselsberg m"usste ich mich dann mit dem begn"ugen, was ich 1821 durch Encke's Heliotrop erhalten habe. Es versteht sich, dass Sie die Lage des neuen Steins und Heliotropplatzes gegen das Haus auf das sch"arfste bestimmen m"ussten. Ich wiederhole jedoch, dass es besser sein w"urde, wenn ich den Winkel vom Inselsberg zum Brocken auf dem Hohehagen unmittelbar erhalten k"onnte.
\item Am allerbesten w"are es wohl, wenn Sie sogleich sich noch einen dritten Heliotrop anschafften. Der kleine Aufwand w"urde gewiss reichlich durch den Genauigkeits- und Zeitgewinn "uberwogen, den Sie k"unftig davon haben w"urden. Dies ginge jetzt um so leichter an, da Rumpf eben einen neuen Heliotrop fertig stehen hat. Er hat nemlich einen f"ur Schumacher verfertigt (und bereits abgesandt) und bei der Gelegenheit, weil, wie er sagt, es sich leichter arbeiten l"asst, sogleich zwei auf einmal gearbeitet. Der zweite ist also in diesem Augenblick noch disponibel, obwohl nicht zu zweifeln ist, dass er auch sonst leicht Gelegenheit finden wird, ihn abzusetzen. (Bei mir sind bereits mehrere Anfragen aus dem Auslande eingegangen.) Ich habe Hrn.\ Rumpf ersucht, diesen Heliotrop theils ganz versendungsfertig zu machen, theils ihn nicht eher an sonst jemand wegzugeben, bis Ihre Erkl"arung darauf erfolgt sein w"urde. Wenn Sie darauf reflectiren, so ist vielleicht das sicherste und k"urzeste, ihn nach dem Meisner schicken oder abholen zu lassen.
\item Ich k"onnte auch etwas fr"uher nach dem Brocken, als Sie nach dem Inselsberge, abgehen; Sie m"ussten auf der Milseburg also suchen, die Richtung zum Inselsberg vor allen andern zu erhalten, damit von dem zu verabredenden Tage an, wo ich auf dem Brocken eintreffe, Ihr Gehilfe vom Inselsberg mir fortw"ahrend Licht zum Brocken schickt. So w"urden Sie Ihren Heliotrop auf dem Inselsberg eine k"urzere Zeit nach Ihrer eigenen Ankunft zu behalten n"othig haben.
\end{enumerate}

\newpage
\noindent Sollte aber alles dieses mit Ihren Pl"anen nicht vereinbar sein, so w"urde ich von der ganzen Expedition, f"ur dies Jahr wenigstens, abstrahiren m"ussen (in welchem Fall ich aber um m"oglichst baldige Benachrichtigung bitte, damit ich und meine Gehilfen uns danach einrichten k"onnen); auch weiss ich nicht, ob ich im k"unftigen Jahre im Stande sein w"urde, mich in demselben Maasse einrichten zu k"onnen wie diesmal. Diesmal nemlich erbiete ich mich, die Zeit des Anfangs ganz Ihnen zu "uberlassen. Zeigen Sie mir nur den Tag genau an, von welchem an Sie mir zuverl"assig Licht vom Inselsberg zum Brocken geben werden. An dem Tage bin ich bestimmt da, und im Fall 1, 2, 3 schicke ich Ihnen auch sogleich Heliotroplicht; im Fall 4, sobald Ihr Licht das einfache Attentionszeichen gibt, nach gew"ohnlichen Uhrschl"agen von $0\frac{3}{4}$, nemlich
\begin{center}
$\displaystyle \frac{0}{|} \quad \frac{0\frac{3}{4}}{|} \quad \frac{1\frac{1}{2}}{|} \quad \frac{1\frac{3}{4}}{|} \quad \frac{2}{|} \quad \frac{2\frac{1}{4}}{|} \quad \frac{3}{|}$ etc.
\end{center}
Doch bemerke ich dabei, dass die Bedeckungen immer recht vollst"andig sein und eher etwas l"anger dauern m"ussen als die "Offnungen; ich meine so:
\begin{center}
$\displaystyle \frac{0}{|} \quad \frac{0\frac{3}{4}}{|} \quad \frac{0\frac{5}{8}}{|} \quad \frac{1\frac{3}{2}}{|} \quad \frac{1\frac{3}{6}}{|}$ etc.
\end{center}
Andere Zeichen gebe ich mit Zahlen, also zum Beispiel die Zahl 3, nach folgendem Schema:

Ich Attentionszeichen von unbestimmter Anzahl, bis Sie dieselben erwiedern und zwar so lange, bis ich aufh"ore, worauf Sie auch aufh"oren und Acht geben. Inzwischen mag etwa $\frac{1}{4}$ Minute gewartet werden. Dann gebe ich das Zeichen selbst und zwar so:
\begin{center}
$\displaystyle \frac{0}{|} \quad \frac{1\frac{2}{5}}{|} \quad \frac{2\frac{9}{0}}{|} \quad \frac{3\frac{8}{6}}{|} \quad \frac{4\frac{8}{0}}{|} \quad \frac{5\frac{1}{2}}{|} \quad \frac{6\frac{0}{0}}{|} \quad \frac{5\frac{8}{8}}{|} \quad \frac{7\frac{8}{8}}{|}$ \\
offen \quad bedeckt \quad offen \quad offen \quad offen \quad offen \quad offen \quad\ldots \\
Nr.~1 \quad\quad Nr.~2 \quad\quad\quad Nr.~3
\end{center}
Die Zahl also $= n$ gesetzt, wechseln

$n+1$ mal bedeckt, 4 Schl"age lang,

mit $n$ mal offen, 1 Schlag lang (quasi >>Blitze<<);

vorher und nachher, wie bestimmt, \emph{infinite} quasi offen.

Nachher wiederholen Sie dasselbe Man"over umgekehrt; nur, insofern nicht lange Unterbrechung stattgefunden hat, ohne meine Erwiederung Ihres Attentionszeichens abzuwarten, sondern nachdem Sie dasselbe etwa 10 bis 20-mal gemacht haben, pausiren Sie (offen) $\frac{1}{4}$ Minute lang und geben das Zeichen.

Doppelte Zeichen, z.~B.\ 3.3, werden durch doppeltes Attentionszeichen angek"undigt ($0\frac{3}{8}$ offen und $0\frac{3}{8}$ zu), und nachdem die erste Zahl gegeben ist, etwa 5 bis $10^{\mathrm{s}}$ lang offen gelassen und dann die zweite Zahl gegeben, also
\begin{center}
$\displaystyle \frac{6\frac{4}{0}}{|} \quad \frac{7\frac{6}{0}}{|} \quad \frac{12\frac{9}{0}}{|} \quad \frac{13\frac{9}{6}}{|} \quad \frac{14\frac{2}{0}}{|} \quad \frac{15\frac{3}{6}}{|} \quad \frac{16\frac{9}{0}}{|} \quad \frac{17\frac{5}{6}}{|} \quad \frac{18\frac{8}{0}}{|} \quad \frac{19\frac{6}{0}}{|}$
\end{center}
wie vorher Nr.~3 \quad $1^{*}$ \quad $2^{*}$ \quad $3^{*}$ \quad offen

\newpage
\noindent Den Tag des Anfangs "uberlasse ich, wie schon gesagt, Ihnen. Nur muss ich so fr"uh davon benachrichtigt werden, dass ich nach Empfang Ihres Briefes noch wenigstens 4 bis 5 Tage habe, um meine Gehilfen in Hannover zu beordern (denn jetzt ruhen die Gesch"afte ganz), damit diese auch zur rechten Zeit am Platze sind.

Noch Einen Umstand muss ich erw"ahnen. Die Gehilfen auf dem Hohehagen und Meisner und eventuelliter, sobald ich auf dem Hohehagen bin, der auf dem Brocken, haben nach 2 Richtungen Licht zu schicken, also \emph{alternative}.

Sie werden das nun \emph{pro aequo} vertheilen, einmal dem einen 1 Stunde oder 2, dann dem andern. Allein es wird gut sein, dass solche jedesmal, wo sie die Absicht haben zu wechseln, dies etwa 5 Minuten vorher durch ein $\frac{1}{2}$ Minute dauerndes Attentionszeichen ank"undigen, damit [d]er [Beobachter] durch das pl"otzliche Abbrechen nicht in einer Beobachtungsreihe im Stiche gelassen zu werden risquire, sondern sie erst geh"orig schliessen k"onne. W"ahrend ich auf dem Hohehagen bin, k"onnte Ihr Gehilfe mir auch wohl vom Meisner aus zuweilen Licht geben. Inzwischen werde ich vermuthlich auch den Stein selbst pointiren k"onnen, sobald ich einmal weiss, wo er steht, wenn er, warum ich bitte, schwarz gef"arbt ist (mit Theer). Doch muss dann der Gehilfe nicht unmittelbar am Stein stehen. Geben Sie ihm also auf, dann, besonders immer wenn die Sonne nicht scheint, und, wenn sie scheint, jedesmal, nachdem der Heliotrop eben eingestellt ist, sich auf ein Dutzend Schritte von der Richtung zum Hohehagen seitw"arts zu halten. Vielleicht geben Sie ihm auch auf, schon fr"uher einmal Licht nach G"ottingen zu schicken. Es soll vom 18.\ d.\ M.\ an t"aglich zwischen 2 und $2\frac{1}{2}$ Uhr aufgepasst und Empfang von Licht durch Erwiederung angezeigt werden. Ich f"urchte aber, dass Ihr Gehilfe M"uhe haben wird, die Richtung heraus zu finden. W"usste ich die Stelle, so w"urde ich das Licht durch eigenes herlocken. Bis jetzt kann ich aber auf dem Gipfel mit dem Teleskop nichts Steinahnliches sehen. Die h"ochste Stelle des Meisner ist hier als ein (im verticalen Sinn) schmaler Saum, und als kahle Bl"osse "uber Holz her sichtbar.

\newpage
\noindent In dem Fall Nr.~4 w"urden Sie Ihre Ankunft auf dem Inselsberg dem Gehilfen auf dem Hohehagen durch einmalige Lichtsendung kund thun. Er liegt $98^{\circ}28'$ links von der Seeberger Sternwarte.

Ich gebe Ihnen anheim, ob Sie auch, wenn ich auf dem Hils bin, mir noch einmal Licht vom Meisner dahin schicken lassen wollen. Es w"urde immer zur vollst"andigen Verbindung beitragen. Herausfordern wollte ich das selber schon, da ich bis dahin den Platz scharf genug kennen werde.

Wenn es m"oglich ist, soll das Heliotroplicht, von dem Platz, wo ich eben selbst messe, im genauen Alignement mit dem Dreieckspunkt sein; geht dies z.~B.\ auf dem Brocken nicht an, so wird nat"urlich die Abweichung genau gemessen.

Am Hercules habe ich vom Hohehagen aus den Kopf pointirt; vom Brocken aus ist der Hercules nur 4-mal geschnitten und auf die ganze Station oder richtiger auf die ganze pyramidalische Spitze des Oktogons pointirt. --- Vom Hohehagen aus ist zwar der Kn"ull unsichtbar; ich vermuthe aber, dass Am"oneburg sichtbar ist, vielleicht auch Hohelohr.\,\ldots\ldots\ldots

\medskip
\noindent \textbf{Gauss an Gerling.} G"ottingen, 1.~September 1823.

\medskip
\noindent Ich habe Ihnen nur den Modus meines Telegraphirens angezeigt; bestimmte bleibende Werthe haben die Zahlzeichen nicht, ebenso wenig wie $a$ und $b$, $y$, $z$ in der Algebra; die Bedeutungen werden immer f"ur solche Umst"ande, als sich eben im voraus erwarten lassen, vorher verabredet. Ich besorge indess, dass das Telegraphiren wenigstens mit dem Meisner-Heliotrop sehr bedenklich sein w"urde, denn es scheint mir, dass er sehr unvollkommen berichtigt sein muss. Besonders in der Stunde von 2--$3\frac{1}{2}$ (wo doch die "ubrigen Gegenst"ande gew"ohnlich zu stark wallen, um etwas brauchbares messen zu k"onnen)\footnote{Nemlich zwischen 6 und $7^{\mathrm{h}}$ habe ich es seltener beachtet, weil ich keine Zeit f"ur Messung verlieren wollte.}, habe ich allezeit die Erfahrung gemacht, dass das Licht, wenn es im besten Leuchten ist, pl"otzlich abgebrochen wird, dann eine Zeit lang ganz unsichtbar bleibt, und dann allm"ahlig ausserst klein anf"angt und zuletzt erst den vollen Glanz erh"alt, aber leider nur kurze Zeit beh"alt (manchmal nur $\frac{1}{4}$ Minute); unter solchen Umst"anden ist aber an zuverl"assiges Telegraphiren nicht zu denken.

\newpage
\noindent Ich habe 1822 einen "ahnlichen Fall mit dem damals auf dem Deister befindlichen "altesten Heliotrop gehabt, dessen Rectification vermuthlich durch unsanften Transport gelitten hatte. Ich half damals durch ein Palliativmittel. Es wurden nemlich "ofters Versuche gemacht, indem ich zuerst dem Lieutenant Hartmann durch ein bestimmtes Zeichen andeutete, dass die Versuche gemacht werden sollten. Dann liess er das Sonnenbild 2 oft wiederholten Malen in verschiedenen H"ohen durch das Fadenkreuz gehen,
\begin{center}
$\uparrow\downarrow$
\end{center}
und ich zeigte ihm den Augenblick, wo ich sein volles Licht zuerst sah, durch Er"offnung meines Heliotrops, und den Schluss des vollen Lichts durch Bedeckung an. Es versteht sich, dass dabei mein Heliotrop immer beinahe central die Sonne auf dem Fadenkreuz hatte, um ganz gewiss zu sein, dass mein Licht v"ollig hinkomme. Meinerseits waren also 3 Personen n"othig: Einer, der immer den Heliotrop fast central unterhielt, ein Zweiter, der mit dem Fernrohr (die Distanz war $8\frac{1}{2}$ Meilen) den jenseitigen Heliotrop beobachtete und ein Dritter, der auf des Zweiten Zurufen den Spiegel des Heliotrops "offnete oder bedeckte (bloss durch Vortreten mit seinem K"orper). Ich wollte dies recht gern auch mit Ihnen thun. Das 3-fache Attentionszeichen ($1\frac{3}{2}$ zu, $1\frac{3}{2}$ offen, $1\frac{3}{2}$ zu, etc.\ etc.) k"onnte die Intention des Versuchs andeuten (den man immer nur macht, wenn wenig Gefahr von Wolkenbedeckung stattfindet), und zwar gibt der das Zeichen zuerst, der den Heliotrop des andern pr"ufen will. Da Gelehrten gut predigen ist, so brauche ich wohl nichts "uber die Benutzung der Versuche beizuf"ugen. Bei "ahnlicher Winkelstellung ist dies Mittel fast so gut wie wirkliche Berichtigung, da derjenige, der den gepr"uften Heliotrop hat, weiss, wie er das Sonnenbild durchgehen lassen muss, und ob er umstellen muss, ehe die Sonne ausgetreten ist, oder nachdem sie schon ein St"uck ausgetreten, etc. Bei dem neuen Heliotrop wird hoffentlich die Berichtigung gut sein.

\newpage
\noindent Was Sie mir von dem Unterschiede Ihres Zeichengebens mit dem meinigen schreiben, ist mir nicht ganz deutlich. Etwas durchaus wesentliches scheint mir zu sein, dass das Attentionszeichen nothwendig erst erwiedert sein muss, ehe das wirkliche Zeichen gegeben wird, ob aber jenes und dieses Zeichen mit einem besondern Fernrohr oder in Ermangelung eines solchen mit dem Heliotropfernrohr beobachtet wird, ist wohl eines und dasselbe.

Der Modus meiner Zeichen ist von mir "ofters ver"andert, derjenige, den ich Ihnen geschrieben, scheint mir der zweckm"assigste. Allein ich habe leider vers"aumt, mir selbst eine Abschrift zu machen, und erinnere schon jetzt bloss aus dem Gedichtniss mich nicht mehr ganz genau. Lassen Sie mir also doch g"utigst diese Stelle meines Briefes copiren.\,\ldots\ldots\ldots

\medskip
\noindent \textbf{Gauss an Gerling.} G"ottingen, 5.~September 1823.

\medskip
\noindent Unbedingt w"unsche ich, dass Sie den Heliotrop [auf dem Inselsberg] genau in das Alignement mit dem Theodolithenplatz stellen und das Zelt so viel wie m"oglich symmetrisch aufschlagen. K"onnten Sie der dem Brocken zugekehrten Seite eine dunkle Farbe geben, so w"are es so viel besser. Ich hoffe, dass der dortige Heliotrop das Zelt weit "uberschreien soll; allein ich habe 1822 doch den Fall gehabt, dass das weisse Zelt in Scharnhorst meine Beobachtungen in Garssen einen ganzen Tag unbrauchbar machte, obwohl dabei zu bemerken, dass erstens das Zelt mir seine S"udseite, von der Sonne beleuchtet, zukehrte, zweitens das Heliotroplicht sehr geschw"acht war, welches beides hier wegfallt. Inzwischen m"ochte ich doch aus Vorsicht den Heliotrop nicht excentrisch aufgestellt wissen, da sonst das Zelt wohl ein wenig beitragen k"onnte, den Winkel zu verf"alschen. Ich vermuthe, wenn z.~B.\ das Zelt $\frac{1}{45}$ so viel Glanz hat wie der Heliotrop, und dieser so weit absteht, dass die Distanz $4''$ gross sein solle, der Winkel um $0\frac{5}{1}$ verf"alscht werden wird.

\medskip
\noindent \textbf{Gauss an Gerling.} G"ottingen, 3.~October 1823.

\medskip
\noindent Erlauben Sie mir noch einige Bemerkungen.
\begin{enumerate}[label=\arabic*.]
\item Beim Vergleichen zweier Punkte fange ich jedesmal von demjenigen an, von dem ich am meisten bef"urchte, dass er versagen k"onnte.
\item Ist der zweite Punkt Heliotroplicht und nicht die allergr"osste Hoffnung, dass er nicht versagen wird, so lese ich "ofters ab als gew"ohnlich. Wo diese Furcht nur ausserst gering war, habe ich wohl zuweilen mehr als 10, d.\ i.\ 15, ja wohl 20 Messungen gemacht, ehe ich wieder ablas. Inzwischen setzt man dabei immer viel aufs Spiel. Ist viel Besorgniss des Versagens des zweiten Punktes, so lese ich wohl nach der dritten, ja nach der zweiten oder gar bei jeder Beobachtung ab. Dasselbe auch, wenn der erste Punkt etwas lange ausbleibt.
\item Bin ich in einer Beobachtungsreihe begriffen, wo der erste Punkt geschnitten ist, so lasse ich gern, ehe ich die Alhidade lese, durch ein Handfernrohr nachsehen, ob der zweite Punkt sichtbar ist; versagt er aber, nachdem schon gelesen ist und kommt nicht wieder, so ist deswegen doch die Messung nicht verloren; ich schneide dann im Nothfall einen dritten Punkt ein und bekomme so \emph{gemischte} Winkel. Bei meinem diesmaligen Brockenbesuch ist indessen dieser Fall nur Einmal eingetreten, indem ich 2-mal (Ilefeld-Hohehagen) + 2-mal (Ilefeld-Hils) nahm. Solche gemischte Winkel werden im System ebenso gut und ganz \emph{pro rata} benutzt wie reine. Im Jahr 1821 habe ich viele gemischte Winkel, sp"ater seltener. Ich habe daher auch an den meisten Stationen, ausser den Hauptrichtungen, Hilfsrichtungen; besonders im Jahre 1821, wo ich nur Einen Heliotrop hatte. Ich w"ahle zu H"ulfsrichtungen gern spitze Th"urme, nicht gar zu entfernt und wo m"oglich nicht gar zu arg ausser der Horizontalebene. Letztere Bedingung war freilich auf dem Brocken nicht zu erreichen. Huyseburg liegt $1^{\circ}37'$, H"uttenrode $1^{\circ}54'$ und eine Stange auf dem Wurmberg, die ich 1821 aufpflanzen und 1823 erneuern liess (obwohl an einem etwas andern Platz), $2^{\circ}13'$ unter dem Horizont.
\end{enumerate}

\medskip
\noindent \textbf{Gauss an Gerling.} G"ottingen, 19.~Julius 1827.

\medskip
\noindent Obgleich ich bei meinen astronomischen Operationen nicht in dem Maasse, wie ich gew"unscht hatte, vom Wetter beg"unstigt bin, so glaube ich doch die Amplitudo des Bogens zwischen den Sternwarten von G"ottingen und Altona bis auf einen sehr kleinen Bruch einer Secunde festgestellt zu haben; ich habe eine sehr grosse Menge von Sternen genommen und zusammen gegen 900 Beobachtungen gemacht. Es bleibt der Unterschied von $5''$ gegen die Rechnung aus den Dreiecken nach Walbeck's Erddimensionen, und es w"are daher sehr interessant gewesen, noch einen Zwischenpunkt zu haben. Mit dem Sector war es unter den obwaltenden Umst"anden unm"oglich; vielleicht aber beobachten wir k"unftig noch mit einem im ersten Vertical aufzustellenden Passageninstrument gemeinschaftlich in Celle, sowie Schumacher vorher in Altona und ich nachher in G"ottingen.\,\ldots\ldots\ldots

\medskip
\noindent \textbf{Gauss an Gerling.} G"ottingen, 12.~September 1838.

\medskip
\noindent Was den L"angenunterschied mit Paris betrifft, so habe ich selbst die Resultate aus Sternbedeckungen etc., die Triesnecker, Wurm u.\ a.\ gezogen haben, niemals zusammengestellt: Sie selbst sind also ebenso gut wie ich selbst im Stande, dieses Resultat zu ermitteln. Dagegen habe ich aus den L"angenunterschieden zwischen Jever und G"ottingen, imgleichen Bentheim und G"ottingen, wie Sie aus meinen, resp.\ meines Sohnes, trigonometrischen Messungen folgten, einerseits, und Jever und Bentheim mit Paris andererseits wie sie Krayenhoff ansetzt, denjenigen L"angenunterschied zwischen G"ottingen und Paris abgeleitet, welchen Sie in H"arpur's Ephemeriden angesetzt finden. Eine neue Rechnung habe ich deshalb nicht gemacht; bei einer completen Ausgleichung aller Dreiecke, die als Kranz das Oldenburgsche umgeben, leidet Jever eine, doch nur sehr geringe Ab"anderung; auch k"onnte Jever ganz weggelassen und dagegen Emden und Onstwedde gebraucht werden, bis wohin

\end{document}
